\documentclass[10pt,letterpaper]{article}
\usepackage[margin=.78in]{geometry}
\usepackage[T1]{fontenc}
\usepackage{lmodern,amsmath,amssymb,amsthm,mathtools}
\usepackage{booktabs,array,longtable,microtype,xcolor,enumitem,fancyhdr}
\usepackage[hidelinks]{hyperref}
\usepackage{xurl}
\definecolor{ink}{HTML}{18334D}
\definecolor{accent}{HTML}{226B70}
\pagestyle{fancy}\fancyhf{}
\fancyhead[L]{\small\sffamily\color{ink}Direct Diophantine consistency}
\fancyhead[R]{\small\sffamily 220 and 218 integer variables}
\fancyfoot[C]{\small\thepage}
\setlength{\headheight}{14pt}
\setlength{\parindent}{0pt}
\setlength{\parskip}{5pt}
\setlength{\emergencystretch}{2em}
\setlist{itemsep=2pt,topsep=4pt}
\newtheorem{theorem}{Theorem}[section]
\newtheorem{lemma}[theorem]{Lemma}
\newtheorem{proposition}[theorem]{Proposition}
\newtheorem{corollary}[theorem]{Corollary}
\theoremstyle{definition}\newtheorem{definition}[theorem]{Definition}
\newtheorem{remark}[theorem]{Remark}
\newcommand{\Z}{\mathbb Z}\newcommand{\Q}{\mathbb Q}
\newcommand{\Con}{\operatorname{Con}}
\newcommand{\Refut}{\operatorname{Ref}}
\newcommand{\Zero}{\mathcal Z}
\newcommand{\htp}{\operatorname{ht}}
\newcommand{\supp}{\operatorname{supp}}
\newcommand{\sha}[1]{\texttt{\scriptsize\detokenize{#1}}}
\hypersetup{pdftitle={Two smaller direct quartics for ZFC consistency},pdfauthor={AI-assisted research prepared for Vladimir Reshetnikov},pdfsubject={Exact integral projections, explicit polynomial artifacts, and a scoped Macaulay rank certificate}}
\begin{document}
\begin{center}
{\sffamily\color{accent}\small PROVEIT RESEARCH NOTE\par}
\vspace{7pt}
{\LARGE\bfseries\color{ink}Two smaller direct quartics\\for ZFC consistency\par}
\vspace{7pt}
{\large Integral projection after pooled power certificates\par}
\vspace{8pt}
Prepared for Vladimir Reshetnikov and the ProveIt research program\\
25 October 2026
\end{center}
\begin{abstract}
We construct two fully materialized integer polynomials whose insolubility
is provably equivalent in ZFC to syntactic consistency of ZFC. The recommended
balanced construction is a quartic in 220 unrestricted integer variables,
with 7,793 collected monomials and coefficient height 45,161. It improves
the inspected direct 225-variable parent simultaneously in arity, number
of quadratic residuals, collected support, and coefficient height. A second
quartic has 218 integer variables, at the cost of more terms and larger
coefficients. Both equations are printed completely in the appendices and
supplied as literal sparse integer data.

The new contribution is an exact composition: five common quadratic-kernel
projections on the pooled-power system expose two further affine relations.
Integral centering and reversible row operations then improve the resulting
representations. We give explicit integer restoration maps, independent
coefficient verification, an authenticated complete parent archive, and a
finite-field certificate excluding affine consequences with affine
multipliers for the 218-variable presentation. The general projection method
is inherited from report 159; the new equations and their measured
composition are the results of this iteration. Global simplicity and
literature-wide priority are not asserted.
\end{abstract}

\section{The equations and the precise result}
\label{sec:result}

All unknowns in this note range independently over \emph{unrestricted}
integers. For the 157 explicitly printed quadratic polynomials
$S_0,\ldots,S_{156}$ in Appendix~\ref{app:220}, define
\begin{equation}\label{eq:p220}
\boxed{\displaystyle P_{220}(x_0,\ldots,x_{219})
       =\sum_{i=0}^{156}S_i(x_0,\ldots,x_{219})^2=0.}
\end{equation}
For the 155 quadratic polynomials $R_0,\ldots,R_{154}$ in
Appendix~\ref{app:218}, define the alternative
\begin{equation}\label{eq:p218}
\boxed{\displaystyle P_{218}(x_0,\ldots,x_{217})
       =\sum_{i=0}^{154}R_i(x_0,\ldots,x_{217})^2=0.}
\end{equation}
These are fixed ordinary polynomials with literal integer coefficients.
The finite sums do not conceal extra unknowns, variable exponents, an
unbounded computation, or a free program parameter. Every declared
coordinate occurs in the corresponding polynomial.

\begin{theorem}[Consistency equations]\label{thm:consistency}
For $n=220$ and $n=218$, the supplied polynomial $P_n$ has exact total
degree four, and
\[
\mathrm{ZFC}\ \vdash\
\bigl(\neg\exists\mathbf x\in\Z^n\ P_n(\mathbf x)=0\bigr)
\ \longleftrightarrow\ \Con(\mathrm{ZFC}).
\]
Here consistency refers to the effective ZFC presentation used by the
identified inherited calibration. The proof composes the ordinary
calibration of the exact $P_{225}$ parent with the integral transformations
proved in Sections~\ref{sec:projection}--\ref{sec:construction}.
\end{theorem}

The logical conclusion is an equivalence theorem, not a proof of either
side. In particular, no assumption that ZFC is consistent is used in
constructing either polynomial. The newly supplied verification consists
of ordinary mathematical proofs and exact finite computations; it is not
a new end-to-end Lean or Rocq formalization.

\subsection{Measured comparison}
\begin{center}\small
\begin{tabular}{lrrr}
\toprule
Measure & Pooled parent $P_{225}$ & $P_{220}$ & $P_{218}$\\
\midrule
Unrestricted integer variables &225&220&218\\
Quadratic residuals &162&157&155\\
Residual monomials &1,177&1,132&1,296\\
Exact polynomial degree &4&4&4\\
Collected polynomial monomials &7,867&7,793&14,410\\
Maximum residual coefficient magnitude &256&128&1,920\\
Maximum polynomial coefficient magnitude &128,446&45,161&5,023,040\\
\bottomrule
\end{tabular}
\end{center}

The balanced equation removes five variables and five rows, removes
74 collected monomials, halves the residual coefficient bound, and reduces
the expanded coefficient height by approximately 64.8\%. The 218-variable
equation removes seven variables and seven rows from the same parent.
Its larger support and height are an explicit tradeoff. Binary arithmetic
operation counts have not been optimized or claimed.

``Simplest'' depends on the measure. The repository already contains
constructions with far fewer variables, including a 37-integer quartic
with a large encoded coefficient tuple and equations of much higher degree
in fewer variables~\cite{collection}. These new equations improve the
specific direct C222 family with moderate literal coefficients. They do
not establish a global minimum or supersede every existing tradeoff.
The distinction between degree, arity, and coefficient size is also central
to work on universal pairs~\cite{bayer2026}.

\textbf{Late comparison check.} During preparation, reports 160--165 were
placed at commit \nolinkurl{897ee2b513fae1e802575abc9f8b6fae896b5014}. We read their
six READMEs and froze that comparison separately; we did not replay those
six packages. Their smallest reported unrestricted-integer direct endpoint
has 234 variables, so none duplicates or dominates either new equation.
Several have fewer monomials or lower coefficient height and remain useful
alternatives. Report 164's 206-variable endpoint is over nonnegative
integers. Report 163 also uses the common-Hessian method, reinforcing the
attribution of the mechanism as prior or parallel work. The full scoped
comparison is retained in \nolinkurl{proofs/LATE_ARRIVALS_160_165.md}.

\subsection{Reproducible mathematical objects}
The ordered sparse residual lists are \texttt{variants/220/equation/equation.json}
and \texttt{equation/equation.json}. The complete collected polynomials are
in the neighboring \texttt{expanded\_polynomial.json} files. An entry
\texttt{[c,[i,j]]} denotes $c x_i x_j$, an empty index list denotes a
constant, and a repeated index denotes a power. All coefficients and indices
are integers; no interpretation of a variable's descriptive name is needed
to evaluate the equation.

For example, the first residual in both final systems is
$-x_2^2+x_3$. The entire ordered list, rather than a fragment or an
existence theorem about a polynomial, fixes each equation. The appendix
exporter reparses every printed mathematical expression and checks its
coefficients against the JSON data.

\section{The authenticated parent and its proof}
\label{sec:parent}

The source comparison is fixed at ProveIt commit
\url{https://github.com/VladimirReshetnikov/ProveIt/commit/953fb4d82d4229f4efc5f9db4cc2bf38ae329bb4}.
At this snapshot the collection extends through report 159; its opening
summary describes an older frontier. Report 158, \emph{Pooling exact power
certificates for Diophantine consistency}, is the immediate parent of the
present construction~\cite{pool158}. Its exact residual-system SHA-256 is
\begin{center}\sha{fee84c408828aba3e7ceebe9a37f58509327063aa7fd078fd49c80adf3af4e56}.\end{center}
The input is 29,598 bytes, with Git blob
\sha{425a70a4960e0c60ebedfcae641f08be15dbdcc0}.
The unmodified original archive is included under \texttt{inherited/}.
Its SHA-256 is
\begin{center}\sha{0e0efbc168fcaac562042778ffcd5bf40f6f74a711ee33915c0ea0750aa10c55}.\end{center}

We use the following identified ordinary theorem from that archive:
\begin{equation}\label{eq:parent-calibration}
\mathrm{ZFC}\vdash
\bigl(\exists\mathbf u\in\Z^{225}\ P_{225}(\mathbf u)=0\bigr)
\longleftrightarrow \Refut(\mathrm{ZFC}).
\end{equation}
This is a substantive imported theorem, not something inferred from a
finite set of numerical tests. Its arithmetic construction is direct from
the C222 proof calculus; it uses neither K215 nor Jones's universal
Diophantine polynomial. The present transformation changes the arithmetic
presentation and preserves its entire integer zero set up to explicit
bijections.

\subsection{Why the pooled power certificate is sound}
The central compression in the parent can be stated cleanly. Fix a nonzero
integer base $n$, an integer Pell center $A$, and canonical signed Pell
pairs $(X_i,Y_i)$ at signed indices $k_i$. Existing coordinates represent
$Z_i=AY_i$. Put
\[
T=2An-n^2-1,\qquad
H=\sum_iY_i(m_i-X_i),\qquad
S=\sum_i\bigl((m_i-X_i+Z_i)^2-Y_i^2\bigr).
\]
Two additional coordinates $K,L$ suffice for the quadratic rows
\begin{align}
K-2H-(2A-n)L&=0,\label{eq:pool1}\\
S-nK+L&=0.\label{eq:pool2}
\end{align}
Elimination gives $S-2nH=TL$. Moreover,
\begin{equation}\label{eq:poolid}
\sum_i\bigl(m_i-X_i+(A-n)Y_i\bigr)^2
   =S-2nH-T\sum_iY_i^2.
\end{equation}
Thus $T$ divides the sum of the power-congruence errors. The existing square
frame is
\[
W=2+\sum_j a_j^2\ge258,\qquad |A|>W^W,
\]
and each actual base, output, and requested index is among its affine
atoms. For $b=\lfloor\sqrt{W-1}\rfloor$, all their absolute values are
at most $b$. Since $T\equiv-1\pmod n$, multiplying the Pell congruences
by $n^b$ clears negative indices. Consequently $T$ divides
\[
E=\sum_i\bigl(m_i n^b-n^{b+k_i}\bigr)^2.
\]
For at most fifteen requests,
\[
0\le E\le4r b^{4b}<W^{2b+1}\le W^W<|T|.
\]
Here $4r\le60<W$, $b^2<W$, $2b+1\le W$, and
$|T|\ge2|A|-W>W^W$. It follows that $E=0$. At nonunit bases this also
excludes negative indices: an integer $m_i$ could not then satisfy
$m_i n^{-k_i}=1$. At bases $1$ and $-1$ the correct signed unit values are
recovered without an unjustified sign claim. Conversely, true powers and
canonical pairs make the required quotients integral and supply $K,L$.

The parent retains three singleton requests for the initial bootstrap.
It interprets the $q$ group even in the temporary unit case $q=1$, excludes
that case using the unchanged divisor identity, and only then invokes
nonunit power recovery at $q+1$. Its strict horizon is existentially chosen
from $16+z^2$ before constructing proof encodings. The linear beta base is
$y=V-1$, with $V=(F^2+1)^t$ and ultimately $F=t!$. Its completeness proof
uses a horizon larger than every actual local parameter, label, and proof
length. These choices are cofinal reconstruction arguments; they are not
assertions that arbitrary earlier witnesses are left unchanged.

\subsection{The logical calibration is an ordinary ZFC proof}
\label{sec:logical}
The arithmetic certificate represents refutability in a square-shell
coding of C222. The inherited shell-recoding proof constructs mutually
inverse recursive constructor-preserving translations to the original
coding and fixes truth $0$ and falsity $1$. It reconstructs the numerical
disequality witness in the guarded bundle after recoding, rather than
assuming that a code translation preserves arithmetic.

The next two precise inherited interfaces are
\begin{quote}\small\ttfamily
KEndpoint.c222\_iff\_zf\_refutation\\
ConstructibleChoice.zf\_refutation\_iff\_zfc\_refutation.
\end{quote}
The selected source declarations are included for inspection. The ordinary
proof uses first-order completeness, tagged and quotient set models,
definable substructures, and constructible models. The relevant
well-founded part is \emph{internal to the structure}: its quantifiers range
over represented sets. An externally well-founded, transitive, or
$\omega$-standard model is not assumed.

These model constructions, their satisfaction predicates, and the involved
sets of finite formulas are available in ZFC. Thus the imported argument
supports~\eqref{eq:parent-calibration} inside ZFC. It should not be described
merely as a primitive-recursive translator between finite refutations:
that would be an additional implementation assertion. The new polynomial
maps below are elementary integer identities and are likewise provable
inside ZFC. Neither step requires an assumption of consistency.

The included audit records the precise reviewed scope. The entire
transitive formal-library closure was not rebuilt, and an encoded
object-language ZFC derivation of the new equivalence was not emitted.
This distinguishes an ordinary mathematical proof of provability from a
new proof-assistant certification.

\section{Integral projections that preserve quadratic degree}
\label{sec:projection}

For an ordered family $F=(F_0,\ldots,F_{m-1})$ of integer polynomials,
write $\Zero_{\Z}(F)$ for its simultaneous integer zero set.

\begin{lemma}[Common Hessian direction]\label{lem:radical}
Let $F_i\in\Z[x_0,\ldots,x_{n-1}]$ have degree at most two. Suppose an
integer vector $d$ satisfies $H_i d=0$ for every Hessian $H_i$, and has a
coordinate $d_p=1$. Let $\ell_i$ be the linear coefficient vector and
$a_i=\ell_i\cdot d$. If $a_r\in\{1,-1\}$ for some row $r$, one can remove
one coordinate and one row, retaining quadratic integer residuals and an
explicit bijection of simultaneous integer zero sets.
\end{lemma}
\begin{proof}
Use coordinates $x=y+td$ with $y_p=0$. The inverse integral coordinate map
is $t=x_p$, $y_i=x_i-d_i x_p$ for $i\ne p$. Taylor expansion, or direct
comparison of coefficients, gives
\[
F_i(y+td)=Q_i(y)+a_i t,\qquad Q_i(y)=F_i(y)|_{y_p=0},
\]
because both mixed and squared terms in $t$ vanish. The pivot equation
forces the unique integer value
\[
t=-a_r Q_r(y).
\]
The remaining residuals are
\begin{equation}\label{eq:project}
G_i(y)=Q_i(y)-a_i a_r Q_r(y)\quad(i\ne r).
\end{equation}
They have degree at most two. A zero of the $G_i$ restores to
$x=y-a_r Q_r(y)d$, an integer zero of all original rows. The inverse
projection just displayed recovers $y$. Conversely, every original zero
has the forced value of $t$ and hence projects to a zero of the $G_i$.
\end{proof}

The condition on every Hessian is stronger than asking that a coordinate
have no square term: all mixed terms must disappear as well. A primitive
direction alone also does not suffice. A nonunit slope could introduce a
divisibility condition. The five actual pivots below all have unit slope,
so that issue does not arise. The general mechanism is credited to
report 159~\cite{projection159}; it is not claimed as a newly discovered
general elimination theorem.

\begin{lemma}[Affine consequence]\label{lem:affine}
Suppose an exact constant integer combination of quadratic residuals is
$\sum_i c_i F_i=hL$, where $h\ne0$, $L$ is affine with integer
coefficients, and some coordinate $x_p$ has coefficient $1$ or $-1$ in
$L$. If $c_r\ne0$, solving $L=0$ for $x_p$, substituting in all remaining
rows, and dropping row $r$ gives a bijection of integer zero sets and
preserves degree at most two.
\end{lemma}
\begin{proof}
The equation $L=0$ has an integer affine restoration formula for $x_p$.
Every old zero satisfies it because $hL=0$. At a new zero, restoration
sets $L=0$ and every retained $F_i=0$. The displayed identity gives
$c_r F_r=0$, hence $F_r=0$ in the integers. Substitution is affine, so
degree cannot increase. The projection and restoration are inverse on
the zero sets.
\end{proof}

An affine elimination can work even when the removed coordinate appears
quadratically in other residuals. It is a different move from
Lemma~\ref{lem:radical}.

\begin{lemma}[Row operations and sum of squares]\label{lem:sos}
Replacing row $F_i$ by $F_i+kF_j$ with $i\ne j$ and $k\in\Z$ preserves
the common zero set; its inverse subtracts $kF_j$. An integral triangular
affine coordinate shear also has an integral inverse. Finally,
\[
\sum_i F_i(\mathbf x)^2=0\quad\Longleftrightarrow\quad
F_i(\mathbf x)=0\ \text{for all }i
\]
for every integer tuple.
\end{lemma}
\begin{proof}
The row and coordinate maps have the stated explicit inverses. Each
integer square is nonnegative, so their finite sum is zero exactly when
each summand is zero.
\end{proof}

\section{The actual reductions}
\label{sec:construction}

\subsection{Five directions: 225 to 220}
The following table is a complete index guide to the five applications
of Lemma~\ref{lem:radical}. Indices are zero based and refer to the system
\emph{immediately before that row of the table}, not to the initial
225-coordinate list.

\begin{center}
\begin{tabular}{crcrr}
\toprule
Input arity&Direction $d$&Pivot $p$&Deleted row $r$&Slope $a_r$\\
\midrule
225&$e_{44}-e_{45}$&44&29&$1$\\
224&$e_{60}-e_{61}-2e_{62}$&60&46&$-1$\\
223&$-e_{62}+e_{64}$&64&49&$1$\\
222&$e_{31}-e_{32}-2e_{33}$&31&16&$1$\\
221&$e_{25}-e_{27}-2e_{28}$&25&10&$1$\\
\bottomrule
\end{tabular}
\end{center}

Each recorded direction annihilates every actual residual Hessian.
Every restoration polynomial is printed as a sparse integer expression
in the certificate. Formula~\eqref{eq:project} specifies the transformed
rows without any high-degree substitution. The independent checker also
performs the full restoration substitution directly, using a different
polynomial representation, and verifies that the resulting identities
have degree at most two.

The resulting 220-variable system has 157 quadratic rows. Applying the
38 elementary integer row operations in the sibling certificate produces
the $S_i$ of Appendix~\ref{app:220}. Their sum of squares is the
recommended balanced equation $P_{220}$.

\subsection{Two additional affine contractions: 220 to 218}
For this branch start with the 220-variable system \emph{before} its
38 final row operations. Its row 81 is exactly
\begin{align*}
0={}&x_{95}-9x_9-9x_{10}-4x_{20}-2x_{23}-2x_{27}-2x_{29}
      -2x_{31}-2x_{33}\\
 &-4x_{36}-2x_{37}-2x_{40}-2x_{43}-4x_{62}-2x_{63}.
\end{align*}
Solve for $x_{95}$ and delete this row. In the resulting 219-variable
system, the exact row combination is
\[
F_{19}-F_6=4+12x_9+16x_{10}+4x_{20}-x_{21}+x_{34}.
\]
The previous deletion does not alter these smaller coordinate indices.
Solve this affine consequence for $x_{34}$ and omit row 6. Its coefficient
in the combination is $-1$, so recoverability is immediate.
Lemma~\ref{lem:affine} gives 218 variables and 155 quadratic rows.

Direct substitution at this point increases coefficients substantially.
The integral triangular shear
\begin{equation}\label{eq:center}
x_{21}=z_{21}+2(1+3z_9+4z_{10}+z_{20}),\qquad x_i=z_i\ (i\ne21)
\end{equation}
centers the paired affine expressions. Its inverse subtracts the same
quantity, which does not involve coordinate 21. Finally apply the
43 elementary integer row operations recorded in the primary certificate.
Renaming the retained $z$ coordinates to $x$ yields the residuals $R_i$
of Appendix~\ref{app:218}.

\begin{theorem}[Exact algebraic comparison]\label{thm:bijection}
The supplied transformations define explicit bijections
\[
\Zero_{\Z}(P_{225})\simeq\Zero_{\Z}(P_{220}),\qquad
\Zero_{\Z}(P_{225})\simeq\Zero_{\Z}(P_{218}).
\]
All coordinate maps and restoration maps are integer polynomial maps on
the relevant zero sets. The two output polynomials have the metrics in
Section~\ref{sec:result}.
\end{theorem}
\begin{proof}
Apply Lemma~\ref{lem:radical} successively using the table. For the balanced
branch apply Lemma~\ref{lem:sos} to its row operations. For the smaller-arity
branch apply Lemma~\ref{lem:affine} twice, the invertible shear
\eqref{eq:center}, and its row operations. Lemma~\ref{lem:sos} converts the
common zero sets into zero sets of the two single polynomials.

All emitted residuals are quadratic, so the degrees are at most four.
Complete exact collection verifies degree four, all monomials, all
coefficient heights, and use of every declared coordinate. For $P_{218}$,
the coefficient of $x_0^4$ is 10, providing a simple explicit nonzero
quartic coefficient. The independently checked full expansions provide
the degree witness for both equations.
\end{proof}

\begin{proof}[Proof of Theorem~\ref{thm:consistency}]
Theorem~\ref{thm:bijection} is an elementary finite algebraic proof
available in ZFC. Compose its zero-existence equivalences with
\eqref{eq:parent-calibration}. Taking negations gives the two stated
insolubility equivalences. This composition uses no consistency assumption.
\end{proof}

\section{A precise stopping certificate for one class of reductions}
\label{sec:rank}

The preceding results do not make the arity 218 optimal. Nevertheless,
one can rigorously identify a class of further moves that cannot work
on this presentation.

Let $V$ be the rational vector space of polynomials of degree at most
three modulo its affine subspace. Form the integer coefficient matrix
whose rows are the images in $V$ of
\[
R_i,\quad x_jR_i
\qquad (0\le i<155,\ 0\le j<218).
\]
Equivalently, retain only coefficients of monomials of degrees two and
three. There are $(218+1)155=33,945$ rows and 137,911 occurring columns.

\begin{proposition}[Affine-multiplier obstruction]\label{prop:rank}
For the literal final $P_{218}$ residual list, this matrix has full row
rank over $\Q$. Its common Hessian kernel is also zero. Therefore no
nonzero family of affine rational polynomials $L_i$ can make
$\sum_i L_iR_i$ affine.
\end{proposition}
\begin{proof}
The supplied checker first validates the canonical residual format and
its exact file identity. It verifies primality of $p=1,000,003$ by trial
division. Gaussian elimination over $\mathbf F_p$ gives rank 218 for the
vertically stacked Hessians and rank 33,945 for the displayed matrix.
Every matrix row is processed, with no truncation or skipped row; the
selected pivot columns and their nonzero leading coefficients are saved.

A full-rank minor modulo $p$ is the reduction of an integer minor whose
determinant is not divisible by $p$. It is therefore nonzero over $\Q$.
If a combination with affine multipliers were affine, its coefficient
vector would give a nonzero rational row dependence in this matrix,
contradicting full row rank. The Hessian conclusion follows in the same way.
\end{proof}

The program \texttt{verify\_rank.py} recomputes these certificates using
only the literal final residuals and standard-library integer arithmetic.
The rank receipt identifies the equation by SHA-256. This obstruction
includes constant-multiplier affine consequences, and goes beyond checking
only the quadratic coefficient matrix.

It excludes neither affine consequences with higher-degree multipliers,
nor new quadratic consequences, nonlinear coordinate changes, integral
radical arguments, or changes to the original proof encoding. It is not a
lower bound for Diophantine equations representing consistency. Exploratory
quadratic-closure calculations were not lifted to characteristic zero and
are not claimed as additional theorems here.

\section{Verification, provenance, and reproduction}
\label{sec:verification}

The construction and independent verifier use different sparse monomial
representations. The builder stores repeated coordinate indices; the
verifier uses exponent vectors and performs the actual restoration
substitutions. It does not import the builder. Their agreement is supported
by the mathematical projection lemmas, rather than treated as a substitute
for those lemmas.

The primary verifier checks 2,426 exact symbolic identities; the
220-variable verifier checks 1,948. Both reject deliberate corruption of
the direction data, an affine or row-operation certificate, a residual
coefficient, and a term of the collected quartic. Normal and optimized
Python receipts agree. Explicit exceptions keep validation active under
\texttt{python3 -O}.

Both certificates use one documented canonical JSON serialization for
the complete equation at every stage. Their stage hashes form continuous
chains, and the builder checks every link. The independent verifier checks
separate mathematical hashes of the arity and ordered residuals, excluding
descriptive metadata. The original parent file additionally has its
authenticated Git blob identity.

The unmodified 225-variable parent was freshly rebuilt in both normal and
optimized Python. All 20 generated JSON files reproduced byte for byte,
with 175 forward and 178 reverse private-wire identities, fourteen Pell
coordinate inverses, and all seventeen power requests checked against the
36-atom frame. Its immediate 245-variable parent was also rebuilt and
audited, including 202 integral monic projections. Historical proof
receipts inside the archive remain historical; the new receipts are named
separately in \texttt{verification/}.

These are exact finite artifact checks. They do not decide whether the
equations have roots, prove consistency, or substitute for the inherited
unbounded arithmetic and logical proofs. Reviews in this package were
independent tasks within an AI-assisted research session, not external
human peer review.

\subsection{Commands}
Python 3.10 or later is sufficient; no third-party Python packages, network
access, or proof assistant are needed for the new finite checks. From the
top-level extracted directory:
\begin{verbatim}
python3 -B build.py
python3 -B verify.py
python3 -B variants/220/build.py
python3 -B variants/220/verify.py
python3 -B verify_rank.py
python3 -B export_equations.py
\end{verbatim}
The builder commands regenerate the supplied equation files from the fixed
source and certificate. The verifier commands audit them independently.
An explicit tuple may be evaluated with
\begin{verbatim}
python3 -B evaluate.py --variant 220 --values tuple.json
\end{verbatim}
The evaluator accepts exactly 220 or 218 JSON integers, as selected; it
rejects booleans, fractional values, and incorrect tuple lengths. It is
an evaluator, not a procedure for deciding consistency.

For the document, run the exporter first and then run twice from
\texttt{report/}:
\begin{verbatim}
pdflatex -no-shell-escape -interaction=nonstopmode \
         -halt-on-error -jobname=ConZFC_Integral_Projection report.tex
\end{verbatim}
The standard TeX Live packages used are listed in the preamble. The PDF
has been rendered and visually inspected. PDF bytes can change under
recompilation; the exact equation bytes are controlled separately.

\section{Further research questions}
\label{sec:future}

\begin{enumerate}
\item \textbf{Use higher-degree consequences with certified recovery.}
Proposition~\ref{prop:rank} resolves the degree-one multiplier search for
this presentation. A degree-two or targeted higher-degree ideal calculation
could reveal a primitive affine consequence. Any successful result must
include an exact polynomial-combination identity and an integral restoration
map; a rational elimination alone is insufficient.

\item \textbf{Turn new quadratic relations into reversible degree-preserving
changes.} Shared Pell multiplication wires give quadratic relations of the
form $Y_iZ_j-Y_jZ_i=0$. Such consequences may permit a better presentation
even when they imply no immediate affine identity. Deleting a defining row
requires a proof that it is recovered at every integer solution, including
zero-coordinate cases.

\item \textbf{Combine structural contraction with semantic changes.}
The current projections occur mainly in syntax-expression wires. A new
proof-certificate syntax or a different batching of powers could reduce the
number of necessary nonlinear wires before algebraic elimination. The
entire compiler, all sign and range guards, and its logical calibration
must be counted.

\item \textbf{Reduce the coefficient cost of 218 variables.}
The two affine contractions raise density and coefficient size. A different
integral coordinate basis, a more effective reversible row basis, or a
combined elimination may lower these costs. The 220-variable equation is
the practical balanced checkpoint while this remains open.

\item \textbf{Formalize the exact new certificates.}
The elementary lemmas and fixed integer identities are suitable for a small
verified checker. A separate end-to-end project should connect the literal
new polynomial to the actual encoded ZFC no-zero sentence and its sealed
consistency sentence. This would close an implementation gap without
changing the ordinary mathematical equivalence proved here.
\end{enumerate}

\begin{thebibliography}{9}
\bibitem{pool158}
ProveIt research collection, report 158,
\emph{Pooling exact power certificates for Diophantine consistency},
25 October 2026. Immutable inspected directory:
\url{https://github.com/VladimirReshetnikov/ProveIt/tree/953fb4d82d4229f4efc5f9db4cc2bf38ae329bb4/Computability/DiophantineConsistency/con-zfc-diophantine-forms/158-pooled-powers}.
The complete authenticated original archive is included in this release.

\bibitem{projection159}
ProveIt research collection, report 159,
\emph{Exact projection lemmas and the scope of the rank certificates},
\texttt{proofs/projection\_scope.md}, with the accompanying projection
constructors and independent audit.
\url{https://github.com/VladimirReshetnikov/ProveIt/tree/953fb4d82d4229f4efc5f9db4cc2bf38ae329bb4/Computability/DiophantineConsistency/con-zfc-diophantine-forms/159-quadratic-projection-234}.

\bibitem{collection}
Vladimir Reshetnikov, ProveIt,
\emph{Con(ZFC) in Diophantine form}, collection README and active equation
metadata, inspected at commit 953fb4d82d4229f4efc5f9db4cc2bf38ae329bb4.
\url{https://github.com/VladimirReshetnikov/ProveIt/tree/953fb4d82d4229f4efc5f9db4cc2bf38ae329bb4/Computability/DiophantineConsistency}.

\bibitem{bayer2026}
Jonas Bayer, Marco David, Malte Ha\ss ler, Yuri Matiyasevich, and Dierk
Schleicher, \emph{Diophantine equations over $\mathbb Z$: universal bounds
and parallel formalization}, arXiv:2506.20909v3, 22 October 2026.
\url{https://arxiv.org/html/2506.20909v3}.
This is background on complexity measures, not a dependency of the new
projection identities.
\end{thebibliography}

\clearpage\appendix
\section{The complete balanced 220-variable equation}
\label{app:220}
The independent coordinates are $x_0,\ldots,x_{219}\in\Z$.
Every row below defines the displayed polynomial $S_i$.
The single equation is $P_{220}=\sum_{i=0}^{156}S_i^2=0$.
These 157 rows contain 1,132 nonzero residual monomials. The complete
collected quartic contains 7,793 monomials and has coefficient height
45,161. No descriptive provenance label imposes an extra condition.
\begingroup\small\setlength{\jot}{2pt}
\setlength{\abovedisplayskip}{5pt}\setlength{\belowdisplayskip}{5pt}
\allowdisplaybreaks[4]
% Generated from the literal JSON. All rows are independently reparsed.
% BEGIN ROW 0
\begin{align*}
S_{0}={}&-x_{2}^{2}+x_{3}
\end{align*}
% END ROW 0
% BEGIN ROW 1
\begin{align*}
S_{1}={}&x_{0}-x_{0}x_{5}-x_{5}+x_{8}
\end{align*}
% END ROW 1
% BEGIN ROW 2
\begin{align*}
S_{2}={}&-x_{2}-x_{5}x_{6}+x_{7}x_{8}
\end{align*}
% END ROW 2
% BEGIN ROW 3
\begin{align*}
S_{3}={}&-1+x_{1}-x_{7}x_{8}
\end{align*}
% END ROW 3
% BEGIN ROW 4
\begin{align*}
S_{4}={}&1+x_{3}-x_{4}x_{19}+x_{17}
\end{align*}
% END ROW 4
% BEGIN ROW 5
\begin{align*}
S_{5}={}&-1+2x_{9}-x_{9}^{2}-2x_{9}x_{10}+2x_{10}-x_{10}^{2}+x_{20}
\end{align*}
% END ROW 5
% BEGIN ROW 6
\begin{align*}
S_{6}={}&-4-12x_{9}-16x_{10}-4x_{20}+x_{21}-x_{34}
\end{align*}
% END ROW 6
% BEGIN ROW 7
\begin{align*}
S_{7}={}&-2+2x_{9}-x_{9}^{2}-4x_{9}x_{21}+4x_{21}-4x_{21}^{2}+x_{22}
\end{align*}
% END ROW 7
% BEGIN ROW 8
\begin{align*}
S_{8}={}&67-60x_{9}-3x_{9}^{2}-60x_{9}x_{21}-24x_{9}x_{22}-116x_{21}-124x_{21}^{2}-64x_{21}x_{22}-31x_{22}-4x_{22}^{2}+x_{23}
\end{align*}
% END ROW 8
% BEGIN ROW 9
\begin{align*}
S_{9}={}&-2-x_{10}^{2}-2x_{10}x_{11}-x_{11}^{2}+x_{24}
\end{align*}
% END ROW 9
% BEGIN ROW 10
\begin{align*}
S_{10}={}&-1+2x_{9}-x_{9}^{2}-2x_{9}x_{11}+2x_{11}-x_{11}^{2}+x_{25}
\end{align*}
% END ROW 10
% BEGIN ROW 11
\begin{align*}
S_{11}={}&15-38x_{9}-29x_{9}^{2}-34x_{9}x_{10}-36x_{9}x_{11}-24x_{9}x_{20}-24x_{9}x_{25}+25x_{10}+38x_{10}^{2}-48x_{10}x_{11}\\
{}&+16x_{10}x_{24}-28x_{10}x_{25}+36x_{11}+70x_{11}^{2}-32x_{11}x_{20}+32x_{11}x_{24}-17x_{20}-8x_{20}x_{25}+40x_{24}-14x_{25}\\
{}&+x_{26}-x_{30}-x_{34}
\end{align*}
% END ROW 11
% BEGIN ROW 12
\begin{align*}
S_{12}={}&-15+14x_{9}-2x_{9}^{2}-8x_{9}x_{10}-8x_{9}x_{24}+4x_{10}^{2}+8x_{10}x_{24}+4x_{24}^{2}-4x_{26}^{2}+x_{27}-3x_{45}
\end{align*}
% END ROW 12
% BEGIN ROW 13
\begin{align*}
S_{13}={}&-6+16x_{11}-16x_{11}^{2}-16x_{11}x_{24}-16x_{11}x_{25}+8x_{24}-4x_{24}^{2}-8x_{24}x_{25}+8x_{25}-4x_{25}^{2}+x_{28}
\end{align*}
% END ROW 13
% BEGIN ROW 14
\begin{align*}
S_{14}={}&-6-48x_{9}x_{11}-48x_{9}x_{25}-24x_{9}x_{28}-64x_{10}x_{11}-64x_{10}x_{25}-32x_{10}x_{28}+7x_{11}+9x_{11}^{2}-16x_{11}x_{20}\\
{}&+8x_{11}x_{25}+4x_{11}x_{28}-16x_{20}x_{25}-8x_{20}x_{28}+2x_{21}+8x_{25}+4x_{28}+x_{29}-3x_{34}-x_{44}
\end{align*}
% END ROW 14
% BEGIN ROW 15
\begin{align*}
S_{15}={}&6-19x_{9}-35x_{9}^{2}-80x_{9}x_{11}+4x_{9}x_{24}-24x_{9}x_{25}+33x_{10}+35x_{10}^{2}+80x_{10}x_{11}+24x_{10}x_{24}\\
{}&-4x_{10}x_{25}+16x_{11}+32x_{11}x_{24}-32x_{11}x_{25}+12x_{24}+4x_{24}^{2}-8x_{25}-4x_{25}^{2}+x_{30}+4x_{30}^{2}-x_{31}
\end{align*}
% END ROW 15
% BEGIN ROW 16
\begin{align*}
S_{16}={}&90-10x_{9}-2x_{9}^{2}-24x_{9}x_{10}-32x_{9}x_{11}-8x_{9}x_{24}-72x_{10}-18x_{10}^{2}-84x_{10}x_{11}-48x_{10}x_{24}-96x_{11}\\
{}&-74x_{11}^{2}-64x_{11}x_{24}-78x_{24}-8x_{24}^{2}-4x_{30}^{2}+x_{31}
\end{align*}
% END ROW 16
% BEGIN ROW 17
\begin{align*}
S_{17}={}&-12+16x_{9}-4x_{9}^{2}-8x_{9}x_{11}-8x_{9}x_{20}-8x_{9}x_{24}+16x_{11}-4x_{11}^{2}-8x_{11}x_{20}-8x_{11}x_{24}+16x_{20}\\
{}&-4x_{20}^{2}-8x_{20}x_{24}+16x_{24}-4x_{24}^{2}+x_{32}
\end{align*}
% END ROW 17
% BEGIN ROW 18
\begin{align*}
S_{18}={}&25-56x_{9}-48x_{9}^{2}-128x_{9}x_{11}-64x_{9}x_{24}-32x_{9}x_{25}-32x_{9}x_{32}-90x_{11}-52x_{11}^{2}-48x_{11}x_{24}\\
{}&-8x_{11}x_{25}-40x_{11}x_{32}-24x_{24}-8x_{24}^{2}-16x_{24}x_{32}-38x_{25}+4x_{25}^{2}-8x_{25}x_{32}-2x_{28}-4x_{32}-4x_{32}^{2}\\
{}&+x_{33}
\end{align*}
% END ROW 18
% BEGIN ROW 19
\begin{align*}
S_{19}={}&8+6x_{9}-8x_{9}^{2}-64x_{9}x_{10}+16x_{9}x_{20}-64x_{10}^{2}+6x_{20}+8x_{20}^{2}-x_{21}+3x_{34}-x_{38}
\end{align*}
% END ROW 19
% BEGIN ROW 20
\begin{align*}
S_{20}={}&-2+2x_{9}-x_{9}^{2}-4x_{9}x_{34}+4x_{34}-4x_{34}^{2}+x_{35}
\end{align*}
% END ROW 20
% BEGIN ROW 21
\begin{align*}
S_{21}={}&67-60x_{9}-3x_{9}^{2}-60x_{9}x_{34}-24x_{9}x_{35}-116x_{34}-124x_{34}^{2}-64x_{34}x_{35}-31x_{35}-4x_{35}^{2}+x_{36}
\end{align*}
% END ROW 21
% BEGIN ROW 22
\begin{align*}
S_{22}={}&-x_{10}^{2}-4x_{10}x_{36}-4x_{36}^{2}+x_{37}
\end{align*}
% END ROW 22
% BEGIN ROW 23
\begin{align*}
S_{23}={}&-1-28x_{9}^{2}-32x_{9}x_{10}-40x_{9}x_{20}+8x_{10}-32x_{10}x_{20}-4x_{20}-12x_{20}^{2}-x_{21}+x_{38}
\end{align*}
% END ROW 23
% BEGIN ROW 24
\begin{align*}
S_{24}={}&49-34x_{10}^{2}-44x_{10}x_{11}-24x_{10}x_{24}-32x_{10}x_{38}+6x_{11}-10x_{11}^{2}-24x_{11}x_{24}-24x_{11}x_{38}-32x_{24}\\
{}&-8x_{24}x_{38}+4x_{38}-4x_{38}^{2}+x_{39}-x_{45}
\end{align*}
% END ROW 24
% BEGIN ROW 25
\begin{align*}
S_{25}={}&8x_{9}-64x_{9}^{2}-96x_{9}x_{11}-32x_{9}x_{25}-32x_{9}x_{39}+6x_{11}-36x_{11}^{2}-24x_{11}x_{25}-24x_{11}x_{39}+2x_{25}\\
{}&-4x_{25}^{2}-8x_{25}x_{39}+4x_{39}-4x_{39}^{2}+x_{40}
\end{align*}
% END ROW 25
% BEGIN ROW 26
\begin{align*}
S_{26}={}&-x_{11}^{2}-2x_{11}x_{12}-x_{12}^{2}+x_{41}
\end{align*}
% END ROW 26
% BEGIN ROW 27
\begin{align*}
S_{27}={}&-11+3x_{9}^{2}+8x_{9}x_{11}-8x_{9}x_{25}-4x_{9}x_{41}+4x_{10}^{2}+8x_{10}x_{24}-8x_{10}x_{25}-8x_{11}+2x_{11}^{2}-8x_{11}x_{41}\\
{}&+4x_{24}^{2}-8x_{24}x_{25}+18x_{25}-4x_{25}^{2}+4x_{26}^{2}-x_{27}+4x_{41}-4x_{41}^{2}+x_{42}
\end{align*}
% END ROW 27
% BEGIN ROW 28
\begin{align*}
S_{28}={}&4x_{9}-18x_{9}^{2}-40x_{9}x_{11}-32x_{9}x_{25}-40x_{9}x_{41}-16x_{9}x_{42}+4x_{11}-24x_{11}^{2}-32x_{11}x_{25}-48x_{11}x_{41}\\
{}&-16x_{11}x_{42}-16x_{25}^{2}-32x_{25}x_{41}-16x_{25}x_{42}+4x_{41}-24x_{41}^{2}-16x_{41}x_{42}-4x_{42}^{2}+x_{43}
\end{align*}
% END ROW 28
% BEGIN ROW 29
\begin{align*}
S_{29}={}&x_{11}-25x_{11}^{2}-40x_{11}x_{25}-20x_{11}x_{28}-16x_{25}^{2}-16x_{25}x_{28}-4x_{28}^{2}+x_{44}
\end{align*}
% END ROW 29
% BEGIN ROW 30
\begin{align*}
S_{30}={}&-1+8x_{10}-4x_{10}^{2}-8x_{10}x_{24}+8x_{24}-4x_{24}^{2}+x_{45}
\end{align*}
% END ROW 30
% BEGIN ROW 31
\begin{align*}
S_{31}={}&-1-16x_{10}^{2}-8x_{10}x_{11}-32x_{10}x_{24}-16x_{10}x_{45}-x_{11}^{2}-8x_{11}x_{24}-4x_{11}x_{45}-16x_{24}^{2}-16x_{24}x_{45}\\
{}&-4x_{45}^{2}+x_{46}
\end{align*}
% END ROW 31
% BEGIN ROW 32
\begin{align*}
S_{32}={}&3+x_{10}-2x_{10}x_{11}-4x_{10}x_{46}-3x_{11}^{2}-8x_{11}x_{46}-x_{24}-4x_{46}^{2}+x_{47}
\end{align*}
% END ROW 32
% BEGIN ROW 33
\begin{align*}
S_{33}={}&-1-4x_{47}-4x_{47}^{2}+x_{48}
\end{align*}
% END ROW 33
% BEGIN ROW 34
\begin{align*}
S_{34}={}&-1-x_{11}^{2}-8x_{11}x_{47}-4x_{11}x_{48}-16x_{47}^{2}-16x_{47}x_{48}-4x_{48}^{2}+x_{49}
\end{align*}
% END ROW 34
% BEGIN ROW 35
\begin{align*}
S_{35}={}&4+x_{10}-2x_{10}x_{11}-4x_{10}x_{49}-3x_{11}^{2}-8x_{11}x_{49}-x_{24}-4x_{49}^{2}+x_{50}
\end{align*}
% END ROW 35
% BEGIN ROW 36
\begin{align*}
S_{36}={}&-4x_{44}^{2}-8x_{44}x_{50}-4x_{50}^{2}+x_{51}
\end{align*}
% END ROW 36
% BEGIN ROW 37
\begin{align*}
S_{37}={}&-9+7x_{9}-x_{9}^{2}-8x_{9}x_{44}-4x_{9}x_{51}+24x_{44}-16x_{44}^{2}-16x_{44}x_{51}+12x_{51}-4x_{51}^{2}+x_{52}
\end{align*}
% END ROW 37
% BEGIN ROW 38
\begin{align*}
S_{38}={}&-1-4x_{52}-4x_{52}^{2}+x_{53}
\end{align*}
% END ROW 38
% BEGIN ROW 39
\begin{align*}
S_{39}={}&-x_{11}^{2}-8x_{11}x_{52}-4x_{11}x_{53}-16x_{52}^{2}-16x_{52}x_{53}-4x_{53}^{2}+x_{54}
\end{align*}
% END ROW 39
% BEGIN ROW 40
\begin{align*}
S_{40}={}&-x_{9}-x_{9}^{2}-4x_{9}x_{11}-4x_{9}x_{54}-4x_{11}-4x_{11}^{2}-8x_{11}x_{54}-4x_{54}-4x_{54}^{2}+x_{55}
\end{align*}
% END ROW 40
% BEGIN ROW 41
\begin{align*}
S_{41}={}&-10-12x_{55}-4x_{55}^{2}+x_{56}
\end{align*}
% END ROW 41
% BEGIN ROW 42
\begin{align*}
S_{42}={}&2-4x_{10}-16x_{10}x_{55}-8x_{10}x_{56}-4x_{24}-16x_{24}x_{55}-8x_{24}x_{56}-x_{45}+8x_{55}-16x_{55}^{2}-16x_{55}x_{56}+4x_{56}\\
{}&-4x_{56}^{2}+x_{57}
\end{align*}
% END ROW 42
% BEGIN ROW 43
\begin{align*}
S_{43}={}&-10+x_{10}-25x_{10}^{2}-40x_{10}x_{24}-20x_{10}x_{57}-16x_{24}^{2}-16x_{24}x_{57}-12x_{50}-4x_{50}^{2}-4x_{57}^{2}+x_{58}
\end{align*}
% END ROW 43
% BEGIN ROW 44
\begin{align*}
S_{44}={}&21+24x_{50}-8x_{50}^{2}-16x_{50}x_{58}-2x_{58}-4x_{58}^{2}+x_{59}
\end{align*}
% END ROW 44
% BEGIN ROW 45
\begin{align*}
S_{45}={}&1+4x_{9}-4x_{9}^{2}-8x_{9}x_{20}-8x_{9}x_{61}+x_{11}-x_{11}^{2}-4x_{11}x_{59}+4x_{20}-4x_{20}^{2}-8x_{20}x_{61}-4x_{59}^{2}+x_{60}\\
{}&+4x_{61}-4x_{61}^{2}
\end{align*}
% END ROW 45
% BEGIN ROW 46
\begin{align*}
S_{46}={}&-4x_{9}-4x_{9}^{2}-4x_{9}x_{11}-8x_{9}x_{20}-x_{11}-x_{11}^{2}-4x_{11}x_{20}-4x_{20}-4x_{20}^{2}+x_{61}
\end{align*}
% END ROW 46
% BEGIN ROW 47
\begin{align*}
S_{47}={}&4x_{9}-24x_{9}^{2}-48x_{9}x_{20}-16x_{9}x_{60}-16x_{9}x_{61}+4x_{20}-24x_{20}^{2}-16x_{20}x_{60}-16x_{20}x_{61}-4x_{60}^{2}\\
{}&+8x_{61}-8x_{61}^{2}+x_{62}
\end{align*}
% END ROW 47
% BEGIN ROW 48
\begin{align*}
S_{48}={}&-2-4x_{62}-4x_{62}^{2}+x_{63}
\end{align*}
% END ROW 48
% BEGIN ROW 49
\begin{align*}
S_{49}={}&-x_{2}x_{14}+x_{64}
\end{align*}
% END ROW 49
% BEGIN ROW 50
\begin{align*}
S_{50}={}&x_{13}-x_{14}-2x_{23}x_{64}-x_{65}
\end{align*}
% END ROW 50
% BEGIN ROW 51
\begin{align*}
S_{51}={}&x_{13}-x_{14}-2x_{27}x_{64}-x_{66}
\end{align*}
% END ROW 51
% BEGIN ROW 52
\begin{align*}
S_{52}={}&-2x_{29}x_{64}+2x_{43}x_{64}-2x_{64}-x_{67}+x_{72}
\end{align*}
% END ROW 52
% BEGIN ROW 53
\begin{align*}
S_{53}={}&x_{13}-x_{14}-2x_{31}x_{64}-x_{68}
\end{align*}
% END ROW 53
% BEGIN ROW 54
\begin{align*}
S_{54}={}&x_{13}-x_{14}-2x_{33}x_{64}+6x_{64}-x_{69}
\end{align*}
% END ROW 54
% BEGIN ROW 55
\begin{align*}
S_{55}={}&x_{13}-x_{14}-4x_{36}x_{64}-2x_{37}x_{64}-x_{70}
\end{align*}
% END ROW 55
% BEGIN ROW 56
\begin{align*}
S_{56}={}&x_{13}-x_{14}-2x_{40}x_{64}-x_{71}
\end{align*}
% END ROW 56
% BEGIN ROW 57
\begin{align*}
S_{57}={}&x_{13}-x_{14}-2x_{43}x_{64}-4x_{64}-x_{72}
\end{align*}
% END ROW 57
% BEGIN ROW 58
\begin{align*}
S_{58}={}&-x_{9}x_{10}+x_{9}x_{11}+x_{10}+x_{10}x_{11}-x_{11}-x_{11}^{2}+x_{73}
\end{align*}
% END ROW 58
% BEGIN ROW 59
\begin{align*}
S_{59}={}&1+2x_{12}+x_{12}^{2}-x_{73}^{2}+x_{74}
\end{align*}
% END ROW 59
% BEGIN ROW 60
\begin{align*}
S_{60}={}&-x_{13}^{2}+2x_{13}x_{14}+2x_{13}x_{76}-x_{14}^{2}-2x_{14}x_{76}-x_{74}^{2}+x_{75}-x_{76}^{2}
\end{align*}
% END ROW 60
% BEGIN ROW 61
\begin{align*}
S_{61}={}&-4x_{62}x_{64}-2x_{63}x_{64}+2x_{64}+x_{76}
\end{align*}
% END ROW 61
% BEGIN ROW 62
\begin{align*}
S_{62}={}&-x_{10}x_{64}-x_{14}+x_{77}x_{78}+x_{79}
\end{align*}
% END ROW 62
% BEGIN ROW 63
\begin{align*}
S_{63}={}&-1+2x_{2}-2x_{2}x_{9}-2x_{2}x_{20}+x_{77}
\end{align*}
% END ROW 63
% BEGIN ROW 64
\begin{align*}
S_{64}={}&-x_{9}x_{215}-x_{13}+x_{78}+x_{215}
\end{align*}
% END ROW 64
% BEGIN ROW 65
\begin{align*}
S_{65}={}&-6x_{9}x_{64}-8x_{10}x_{64}+x_{10}x_{215}+x_{13}-x_{14}-2x_{20}x_{64}+x_{80}
\end{align*}
% END ROW 65
% BEGIN ROW 66
\begin{align*}
S_{66}={}&-x_{13}x_{65}+x_{14}x_{65}+x_{81}
\end{align*}
% END ROW 66
% BEGIN ROW 67
\begin{align*}
S_{67}={}&-x_{66}x_{81}+x_{82}
\end{align*}
% END ROW 67
% BEGIN ROW 68
\begin{align*}
S_{68}={}&-x_{67}x_{82}+x_{83}
\end{align*}
% END ROW 68
% BEGIN ROW 69
\begin{align*}
S_{69}={}&-x_{68}x_{83}+x_{84}
\end{align*}
% END ROW 69
% BEGIN ROW 70
\begin{align*}
S_{70}={}&-x_{69}x_{84}+x_{85}
\end{align*}
% END ROW 70
% BEGIN ROW 71
\begin{align*}
S_{71}={}&-x_{70}x_{85}+x_{86}
\end{align*}
% END ROW 71
% BEGIN ROW 72
\begin{align*}
S_{72}={}&-x_{71}x_{86}+x_{87}
\end{align*}
% END ROW 72
% BEGIN ROW 73
\begin{align*}
S_{73}={}&-x_{72}x_{87}+x_{88}
\end{align*}
% END ROW 73
% BEGIN ROW 74
\begin{align*}
S_{74}={}&-x_{79}x_{88}+x_{89}
\end{align*}
% END ROW 74
% BEGIN ROW 75
\begin{align*}
S_{75}={}&-x_{80}x_{89}+x_{90}
\end{align*}
% END ROW 75
% BEGIN ROW 76
\begin{align*}
S_{76}={}&-x_{90}^{2}+x_{91}
\end{align*}
% END ROW 76
% BEGIN ROW 77
\begin{align*}
S_{77}={}&-x_{5}-x_{5}x_{17}+x_{17}+x_{92}
\end{align*}
% END ROW 77
% BEGIN ROW 78
\begin{align*}
S_{78}={}&-x_{15}x_{92}+x_{93}
\end{align*}
% END ROW 78
% BEGIN ROW 79
\begin{align*}
S_{79}={}&-2x_{1}^{2}+2x_{1}x_{13}+2x_{1}x_{14}+2x_{1}x_{93}+2x_{1}x_{94}-x_{13}^{2}-2x_{13}x_{93}-x_{14}^{2}-2x_{14}x_{94}+x_{18}x_{96}-x_{75}x_{91}\\
{}&-x_{93}^{2}-x_{94}^{2}
\end{align*}
% END ROW 79
% BEGIN ROW 80
\begin{align*}
S_{80}={}&-x_{5}x_{16}+x_{16}-x_{16}x_{92}+x_{94}
\end{align*}
% END ROW 80
% BEGIN ROW 81
\begin{align*}
S_{81}={}&-9x_{9}-9x_{10}-4x_{20}-2x_{23}-2x_{27}-2x_{29}-2x_{31}-2x_{33}-4x_{36}-2x_{37}-2x_{40}-2x_{43}-4x_{62}-2x_{63}\\
{}&+x_{95}
\end{align*}
% END ROW 81
% BEGIN ROW 82
\begin{align*}
S_{82}={}&-x_{97}x_{100}+x_{100}+x_{102}
\end{align*}
% END ROW 82
% BEGIN ROW 83
\begin{align*}
S_{83}={}&-x_{4}x_{123}+x_{98}x_{205}-x_{101}-x_{102}x_{123}
\end{align*}
% END ROW 83
% BEGIN ROW 84
\begin{align*}
S_{84}={}&x_{103}x_{205}-x_{104}x_{126}-x_{105}
\end{align*}
% END ROW 84
% BEGIN ROW 85
\begin{align*}
S_{85}={}&x_{106}x_{205}-x_{107}x_{126}-x_{108}
\end{align*}
% END ROW 85
% BEGIN ROW 86
\begin{align*}
S_{86}={}&x_{109}x_{205}-x_{110}x_{126}-x_{111}
\end{align*}
% END ROW 86
% BEGIN ROW 87
\begin{align*}
S_{87}={}&x_{112}x_{205}-x_{113}x_{126}-x_{114}
\end{align*}
% END ROW 87
% BEGIN ROW 88
\begin{align*}
S_{88}={}&x_{115}x_{206}-x_{117}x_{126}-x_{118}
\end{align*}
% END ROW 88
% BEGIN ROW 89
\begin{align*}
S_{89}={}&x_{119}x_{206}-x_{120}x_{126}-x_{121}
\end{align*}
% END ROW 89
% BEGIN ROW 90
\begin{align*}
S_{90}={}&-x_{18}x_{123}+x_{126}
\end{align*}
% END ROW 90
% BEGIN ROW 91
\begin{align*}
S_{91}={}&x_{122}-x_{124}x_{126}-x_{125}
\end{align*}
% END ROW 91
% BEGIN ROW 92
\begin{align*}
S_{92}={}&-x_{126}x_{128}+x_{127}-x_{129}
\end{align*}
% END ROW 92
% BEGIN ROW 93
\begin{align*}
S_{93}={}&-x_{126}x_{131}+x_{130}x_{205}-x_{132}
\end{align*}
% END ROW 93
% BEGIN ROW 94
\begin{align*}
S_{94}={}&-x_{126}x_{134}+x_{133}x_{205}-x_{135}
\end{align*}
% END ROW 94
% BEGIN ROW 95
\begin{align*}
S_{95}={}&x_{136}-2x_{136}x_{138}+x_{139}
\end{align*}
% END ROW 95
% BEGIN ROW 96
\begin{align*}
S_{96}={}&-x_{136}x_{139}-x_{139}+x_{140}
\end{align*}
% END ROW 96
% BEGIN ROW 97
\begin{align*}
S_{97}={}&2x_{137}+x_{137}^{2}+x_{139}^{2}-x_{140}^{2}
\end{align*}
% END ROW 97
% BEGIN ROW 98
\begin{align*}
S_{98}={}&-x_{137}x_{142}-x_{142}+x_{143}
\end{align*}
% END ROW 98
% BEGIN ROW 99
\begin{align*}
S_{99}={}&-1+x_{141}^{2}+2x_{141}x_{143}+x_{142}^{2}
\end{align*}
% END ROW 99
% BEGIN ROW 100
\begin{align*}
S_{100}={}&x_{3}^{2}-2x_{3}x_{137}-2x_{137}+x_{144}
\end{align*}
% END ROW 100
% BEGIN ROW 101
\begin{align*}
S_{101}={}&-x_{3}x_{142}+x_{5}-x_{141}-x_{142}+x_{144}x_{145}
\end{align*}
% END ROW 101
% BEGIN ROW 102
\begin{align*}
S_{102}={}&1-2x_{4}+x_{4}^{2}-2x_{4}x_{137}+x_{146}
\end{align*}
% END ROW 102
% BEGIN ROW 103
\begin{align*}
S_{103}={}&-x_{4}x_{178}+x_{18}+x_{146}x_{147}-x_{177}
\end{align*}
% END ROW 103
% BEGIN ROW 104
\begin{align*}
S_{104}={}&-1+2x_{137}x_{189}-x_{148}+2x_{189}-x_{189}^{2}
\end{align*}
% END ROW 104
% BEGIN ROW 105
\begin{align*}
S_{105}={}&-1+x_{97}+x_{148}x_{149}-x_{177}-x_{178}x_{189}
\end{align*}
% END ROW 105
% BEGIN ROW 106
\begin{align*}
S_{106}={}&-x_{137}x_{151}-x_{151}+x_{152}
\end{align*}
% END ROW 106
% BEGIN ROW 107
\begin{align*}
S_{107}={}&-1+x_{150}^{2}+2x_{150}x_{152}+x_{151}^{2}
\end{align*}
% END ROW 107
% BEGIN ROW 108
\begin{align*}
S_{108}={}&1+x_{137}-x_{137}x_{154}-x_{154}+x_{155}
\end{align*}
% END ROW 108
% BEGIN ROW 109
\begin{align*}
S_{109}={}&x_{153}^{2}+2x_{153}x_{155}-2x_{154}+x_{154}^{2}
\end{align*}
% END ROW 109
% BEGIN ROW 110
\begin{align*}
S_{110}={}&-x_{137}x_{157}-x_{157}+x_{158}
\end{align*}
% END ROW 110
% BEGIN ROW 111
\begin{align*}
S_{111}={}&-1+x_{156}^{2}+2x_{156}x_{158}+x_{157}^{2}
\end{align*}
% END ROW 111
% BEGIN ROW 112
\begin{align*}
S_{112}={}&-x_{137}x_{160}-x_{160}+x_{161}
\end{align*}
% END ROW 112
% BEGIN ROW 113
\begin{align*}
S_{113}={}&-1+x_{159}^{2}+2x_{159}x_{161}+x_{160}^{2}
\end{align*}
% END ROW 113
% BEGIN ROW 114
\begin{align*}
S_{114}={}&-x_{137}x_{163}-x_{163}+x_{164}
\end{align*}
% END ROW 114
% BEGIN ROW 115
\begin{align*}
S_{115}={}&-1+x_{162}^{2}+2x_{162}x_{164}+x_{163}^{2}
\end{align*}
% END ROW 115
% BEGIN ROW 116
\begin{align*}
S_{116}={}&-x_{137}x_{166}-x_{166}+x_{167}
\end{align*}
% END ROW 116
% BEGIN ROW 117
\begin{align*}
S_{117}={}&-1+x_{165}^{2}+2x_{165}x_{167}+x_{166}^{2}
\end{align*}
% END ROW 117
% BEGIN ROW 118
\begin{align*}
S_{118}={}&-x_{137}x_{169}-x_{169}+x_{170}
\end{align*}
% END ROW 118
% BEGIN ROW 119
\begin{align*}
S_{119}={}&-1+x_{168}^{2}+2x_{168}x_{170}+x_{169}^{2}
\end{align*}
% END ROW 119
% BEGIN ROW 120
\begin{align*}
S_{120}={}&-x_{137}x_{172}-x_{172}+x_{173}
\end{align*}
% END ROW 120
% BEGIN ROW 121
\begin{align*}
S_{121}={}&-1+x_{171}^{2}+2x_{171}x_{173}+x_{172}^{2}
\end{align*}
% END ROW 121
% BEGIN ROW 122
\begin{align*}
S_{122}={}&-x_{137}x_{175}-x_{175}+x_{176}
\end{align*}
% END ROW 122
% BEGIN ROW 123
\begin{align*}
S_{123}={}&-1+x_{174}^{2}+2x_{174}x_{176}+x_{175}^{2}
\end{align*}
% END ROW 123
% BEGIN ROW 124
\begin{align*}
S_{124}={}&1-x_{137}x_{178}-x_{178}+x_{179}
\end{align*}
% END ROW 124
% BEGIN ROW 125
\begin{align*}
S_{125}={}&-1+2x_{177}+x_{177}^{2}+2x_{177}x_{179}+x_{178}^{2}
\end{align*}
% END ROW 125
% BEGIN ROW 126
\begin{align*}
S_{126}={}&-x_{137}x_{181}-x_{181}+x_{182}
\end{align*}
% END ROW 126
% BEGIN ROW 127
\begin{align*}
S_{127}={}&-1+x_{180}^{2}+2x_{180}x_{182}+x_{181}^{2}
\end{align*}
% END ROW 127
% BEGIN ROW 128
\begin{align*}
S_{128}={}&-x_{137}x_{184}-x_{184}+x_{185}
\end{align*}
% END ROW 128
% BEGIN ROW 129
\begin{align*}
S_{129}={}&-1+x_{183}^{2}+2x_{183}x_{185}+x_{184}^{2}
\end{align*}
% END ROW 129
% BEGIN ROW 130
\begin{align*}
S_{130}={}&-x_{137}x_{187}-x_{187}+x_{188}
\end{align*}
% END ROW 130
% BEGIN ROW 131
\begin{align*}
S_{131}={}&-1+x_{186}^{2}+2x_{186}x_{188}+x_{187}^{2}
\end{align*}
% END ROW 131
% BEGIN ROW 132
\begin{align*}
S_{132}={}&-4-2x_{0}-3x_{0}^{2}-2x_{0}x_{3}-x_{1}^{2}-2x_{3}-2x_{3}^{2}-x_{4}^{2}-2x_{5}^{2}+2x_{9}-x_{9}^{2}+2x_{10}x_{11}-x_{12}^{2}-x_{13}^{2}\\
{}&-x_{14}^{2}-x_{15}^{2}-x_{16}^{2}-x_{17}^{2}-x_{18}^{2}-x_{24}-x_{95}^{2}+2x_{97}-2x_{97}^{2}-x_{98}^{2}-x_{99}^{2}-x_{103}^{2}-x_{106}^{2}-x_{109}^{2}\\
{}&-x_{112}^{2}-x_{115}^{2}-x_{116}^{2}-x_{119}^{2}-x_{122}^{2}-x_{123}^{2}-x_{127}^{2}-x_{130}^{2}-x_{133}^{2}+x_{136}-x_{189}^{2}
\end{align*}
% END ROW 132
% BEGIN ROW 133
\begin{align*}
S_{133}={}&x_{4}-2x_{4}x_{9}-2x_{4}x_{10}-2x_{4}x_{11}-2x_{4}x_{12}-2x_{4}x_{13}-2x_{4}x_{14}-2x_{4}x_{15}-2x_{4}x_{16}-2x_{4}x_{17}+x_{189}
\end{align*}
% END ROW 133
% BEGIN ROW 134
\begin{align*}
S_{134}={}&-16+x_{0}-x_{190}^{2}
\end{align*}
% END ROW 134
% BEGIN ROW 135
\begin{align*}
S_{135}={}&-x_{137}x_{192}-x_{192}+x_{201}
\end{align*}
% END ROW 135
% BEGIN ROW 136
\begin{align*}
S_{136}={}&-1+x_{191}^{2}+x_{192}^{2}-x_{201}^{2}
\end{align*}
% END ROW 136
% BEGIN ROW 137
\begin{align*}
S_{137}={}&-x_{194}-x_{194}x_{195}+x_{202}
\end{align*}
% END ROW 137
% BEGIN ROW 138
\begin{align*}
S_{138}={}&-1+x_{193}^{2}+x_{194}^{2}-x_{202}^{2}
\end{align*}
% END ROW 138
% BEGIN ROW 139
\begin{align*}
S_{139}={}&-x_{187}^{2}+x_{203}
\end{align*}
% END ROW 139
% BEGIN ROW 140
\begin{align*}
S_{140}={}&x_{192}-2x_{197}x_{203}-x_{203}
\end{align*}
% END ROW 140
% BEGIN ROW 141
\begin{align*}
S_{141}={}&-4x_{192}-8x_{192}x_{196}+x_{195}
\end{align*}
% END ROW 141
% BEGIN ROW 142
\begin{align*}
S_{142}={}&-x_{137}-x_{191}x_{198}+x_{195}
\end{align*}
% END ROW 142
% BEGIN ROW 143
\begin{align*}
S_{143}={}&-x_{186}-x_{188}-x_{191}x_{199}+x_{193}
\end{align*}
% END ROW 143
% BEGIN ROW 144
\begin{align*}
S_{144}={}&-x_{95}-4x_{187}x_{200}+x_{194}
\end{align*}
% END ROW 144
% BEGIN ROW 145
\begin{align*}
S_{145}={}&26+2x_{0}^{2}+2x_{0}x_{3}-2x_{0}x_{142}-2x_{0}x_{151}+x_{1}^{2}-2x_{1}x_{178}+x_{3}^{2}-2x_{3}x_{151}+x_{5}^{2}-2x_{5}x_{169}+x_{9}^{2}\\
{}&-2x_{9}x_{154}+x_{10}^{2}-2x_{10}x_{157}+x_{11}^{2}-2x_{11}x_{160}+x_{12}^{2}-2x_{12}x_{163}+x_{13}^{2}-2x_{13}x_{166}+x_{14}^{2}-2x_{14}x_{172}\\
{}&+x_{15}^{2}-2x_{15}x_{175}+x_{16}^{2}-2x_{16}x_{181}+x_{17}^{2}-2x_{17}x_{184}+x_{95}^{2}-2x_{95}x_{187}-x_{137}x_{204}-2x_{141}-2x_{143}\\
{}&+x_{143}^{2}-2x_{150}-2x_{152}+x_{152}^{2}-2x_{153}+2x_{154}-2x_{155}+x_{155}^{2}-2x_{156}-2x_{158}+x_{158}^{2}-2x_{159}-2x_{161}\\
{}&+x_{161}^{2}-2x_{162}-2x_{164}+x_{164}^{2}-2x_{165}-2x_{167}+x_{167}^{2}-2x_{168}-2x_{170}+x_{170}^{2}-2x_{171}-2x_{173}+x_{173}^{2}\\
{}&-2x_{174}-2x_{176}+x_{176}^{2}-2x_{177}+x_{179}^{2}-2x_{180}-2x_{182}+x_{182}^{2}-2x_{183}-2x_{185}+x_{185}^{2}-2x_{186}-2x_{188}\\
{}&+x_{188}^{2}
\end{align*}
% END ROW 145
% BEGIN ROW 146
\begin{align*}
S_{146}={}&-x_{99}x_{205}+x_{123}
\end{align*}
% END ROW 146
% BEGIN ROW 147
\begin{align*}
S_{147}={}&-x_{116}x_{206}+x_{123}
\end{align*}
% END ROW 147
% BEGIN ROW 148
\begin{align*}
S_{148}={}&-2-x_{95}^{2}-x_{142}^{2}-x_{151}^{2}+2x_{154}-x_{154}^{2}-x_{157}^{2}-x_{160}^{2}-x_{163}^{2}-x_{166}^{2}-x_{169}^{2}-x_{172}^{2}-x_{175}^{2}-x_{178}^{2}\\
{}&-x_{181}^{2}-x_{184}^{2}+x_{187}^{2}-x_{207}^{2}-x_{208}^{2}-x_{209}-x_{209}^{2}
\end{align*}
% END ROW 148
% BEGIN ROW 149
\begin{align*}
S_{149}={}&-1-4x_{101}^{2}-4x_{105}^{2}-4x_{108}^{2}-4x_{111}^{2}-4x_{114}^{2}-4x_{118}^{2}-4x_{121}^{2}+x_{123}^{2}-4x_{125}^{2}-4x_{129}^{2}-4x_{132}^{2}\\
{}&-4x_{135}^{2}-x_{210}^{2}-x_{211}^{2}-x_{212}-x_{212}^{2}-4x_{214}^{2}
\end{align*}
% END ROW 149
% BEGIN ROW 150
\begin{align*}
S_{150}={}&-x_{4}x_{213}+x_{123}
\end{align*}
% END ROW 150
% BEGIN ROW 151
\begin{align*}
S_{151}={}&x_{2}x_{123}-x_{5}x_{213}-x_{214}
\end{align*}
% END ROW 151
% BEGIN ROW 152
\begin{align*}
S_{152}={}&-x_{2}x_{13}+x_{215}
\end{align*}
% END ROW 152
% BEGIN ROW 153
\begin{align*}
S_{153}={}&x_{97}x_{217}-2x_{98}x_{151}+2x_{103}-2x_{103}x_{154}-2x_{106}x_{157}-2x_{109}x_{160}-2x_{112}x_{163}-2x_{115}x_{166}-2x_{119}x_{172}\\
{}&-2x_{122}x_{175}-2x_{127}x_{181}-2x_{130}x_{184}-2x_{133}x_{187}-2x_{137}x_{217}+2x_{150}x_{151}+2x_{151}x_{152}-2x_{153}+2x_{153}x_{154}\\
{}&+2x_{154}x_{155}-2x_{155}+2x_{156}x_{157}+2x_{157}x_{158}+2x_{159}x_{160}+2x_{160}x_{161}+2x_{162}x_{163}+2x_{163}x_{164}+2x_{165}x_{166}\\
{}&+2x_{166}x_{167}+2x_{171}x_{172}+2x_{172}x_{173}+2x_{174}x_{175}+2x_{175}x_{176}+2x_{180}x_{181}+2x_{181}x_{182}+2x_{183}x_{184}+2x_{184}x_{185}\\
{}&+2x_{186}x_{187}+2x_{187}x_{188}+x_{216}-2x_{217}
\end{align*}
% END ROW 153
% BEGIN ROW 154
\begin{align*}
S_{154}={}&1+x_{95}^{2}-x_{97}x_{216}+x_{98}^{2}-2x_{98}x_{150}+x_{103}^{2}-2x_{103}x_{153}+x_{106}^{2}-2x_{106}x_{156}+x_{109}^{2}-2x_{109}x_{159}+x_{112}^{2}\\
{}&-2x_{112}x_{162}+x_{115}^{2}-2x_{115}x_{165}+x_{119}^{2}-2x_{119}x_{171}+x_{122}^{2}-2x_{122}x_{174}+x_{127}^{2}-2x_{127}x_{180}+x_{130}^{2}\\
{}&-2x_{130}x_{183}+x_{133}^{2}-2x_{133}x_{186}+x_{142}^{2}+x_{150}^{2}+x_{153}^{2}+x_{156}^{2}+x_{159}^{2}+x_{162}^{2}+x_{165}^{2}+x_{169}^{2}+x_{171}^{2}\\
{}&+x_{174}^{2}+x_{178}^{2}+x_{180}^{2}+x_{183}^{2}+x_{186}^{2}-2x_{187}^{2}+x_{207}^{2}+x_{208}^{2}+x_{209}+x_{209}^{2}+x_{217}
\end{align*}
% END ROW 154
% BEGIN ROW 155
\begin{align*}
S_{155}={}&x_{97}x_{219}-2x_{99}x_{142}-2x_{116}x_{169}-2x_{123}x_{178}-2x_{137}x_{219}+2x_{141}x_{142}+2x_{142}x_{143}+2x_{168}x_{169}+2x_{169}x_{170}\\
{}&+2x_{177}x_{178}+2x_{178}+2x_{178}x_{179}+x_{218}-3x_{219}
\end{align*}
% END ROW 155
% BEGIN ROW 156
\begin{align*}
S_{156}={}&-x_{97}x_{218}+x_{99}^{2}-2x_{99}x_{141}+x_{116}^{2}-2x_{116}x_{168}+x_{123}^{2}-2x_{123}x_{177}+x_{141}^{2}-x_{142}^{2}+x_{168}^{2}-x_{169}^{2}+x_{177}^{2}\\
{}&-x_{178}^{2}+x_{218}+x_{219}
\end{align*}
% END ROW 156

\endgroup

\clearpage
\section{The complete 218-variable alternative}
\label{app:218}
The independent coordinates are $x_0,\ldots,x_{217}\in\Z$.
The single equation is $P_{218}=\sum_{i=0}^{154}R_i^2=0$.
These 155 rows contain 1,296 nonzero residual monomials. The complete
collected quartic contains 14,410 monomials and has coefficient height
5,023,040. The indexing is specific to this branch and should not be
identified coordinate by coordinate with the preceding appendix.
\begingroup\small\setlength{\jot}{2pt}
\setlength{\abovedisplayskip}{5pt}\setlength{\belowdisplayskip}{5pt}
\allowdisplaybreaks[4]
% Generated from the literal JSON. All rows are independently reparsed.
% BEGIN ROW 0
\begin{align*}
R_{0}={}&-x_{2}^{2}+x_{3}
\end{align*}
% END ROW 0
% BEGIN ROW 1
\begin{align*}
R_{1}={}&x_{0}-x_{0}x_{5}-x_{5}+x_{8}
\end{align*}
% END ROW 1
% BEGIN ROW 2
\begin{align*}
R_{2}={}&-x_{2}-x_{5}x_{6}+x_{7}x_{8}
\end{align*}
% END ROW 2
% BEGIN ROW 3
\begin{align*}
R_{3}={}&-1+x_{1}-x_{7}x_{8}
\end{align*}
% END ROW 3
% BEGIN ROW 4
\begin{align*}
R_{4}={}&1+x_{3}-x_{4}x_{19}+x_{17}
\end{align*}
% END ROW 4
% BEGIN ROW 5
\begin{align*}
R_{5}={}&-1+2x_{9}-x_{9}^{2}-2x_{9}x_{10}+2x_{10}-x_{10}^{2}+x_{20}
\end{align*}
% END ROW 5
% BEGIN ROW 6
\begin{align*}
R_{6}={}&-12+4x_{9}-2x_{9}^{2}-8x_{9}x_{21}-8x_{21}^{2}+x_{22}+x_{34}
\end{align*}
% END ROW 6
% BEGIN ROW 7
\begin{align*}
R_{7}={}&-351-1434x_{9}+412x_{9}^{2}+512x_{9}x_{10}+128x_{9}x_{20}+64x_{9}x_{21}-408x_{9}x_{22}-1920x_{10}-512x_{10}x_{22}\\
{}&-480x_{20}-128x_{20}x_{22}-240x_{21}-64x_{21}x_{22}-190x_{22}-4x_{22}^{2}+x_{23}
\end{align*}
% END ROW 7
% BEGIN ROW 8
\begin{align*}
R_{8}={}&-2-x_{10}^{2}-2x_{10}x_{11}-x_{11}^{2}+x_{24}
\end{align*}
% END ROW 8
% BEGIN ROW 9
\begin{align*}
R_{9}={}&-1+2x_{9}-x_{9}^{2}-2x_{9}x_{11}+2x_{11}-x_{11}^{2}+x_{25}
\end{align*}
% END ROW 9
% BEGIN ROW 10
\begin{align*}
R_{10}={}&17-36x_{9}-27x_{9}^{2}-30x_{9}x_{10}-36x_{9}x_{11}-24x_{9}x_{20}-24x_{9}x_{25}+29x_{10}+39x_{10}^{2}-50x_{10}x_{11}\\
{}&+16x_{10}x_{24}-28x_{10}x_{25}+36x_{11}+69x_{11}^{2}-32x_{11}x_{20}+32x_{11}x_{24}-17x_{20}-8x_{20}x_{25}-x_{21}+41x_{24}\\
{}&-14x_{25}+x_{26}-x_{30}
\end{align*}
% END ROW 10
% BEGIN ROW 11
\begin{align*}
R_{11}={}&-15+14x_{9}-2x_{9}^{2}-8x_{9}x_{10}-8x_{9}x_{24}+4x_{10}^{2}+8x_{10}x_{24}+4x_{24}^{2}-4x_{26}^{2}+x_{27}-3x_{44}
\end{align*}
% END ROW 11
% BEGIN ROW 12
\begin{align*}
R_{12}={}&-6+16x_{11}-16x_{11}^{2}-16x_{11}x_{24}-16x_{11}x_{25}+8x_{24}-4x_{24}^{2}-8x_{24}x_{25}+8x_{25}-4x_{25}^{2}+x_{28}
\end{align*}
% END ROW 12
% BEGIN ROW 13
\begin{align*}
R_{13}={}&4+30x_{9}-48x_{9}x_{11}-48x_{9}x_{25}-24x_{9}x_{28}+40x_{10}-64x_{10}x_{11}-64x_{10}x_{25}-32x_{10}x_{28}+7x_{11}\\
{}&+9x_{11}^{2}-16x_{11}x_{20}+8x_{11}x_{25}+4x_{11}x_{28}+10x_{20}-16x_{20}x_{25}-8x_{20}x_{28}-x_{21}+8x_{25}+4x_{28}+x_{29}-x_{43}
\end{align*}
% END ROW 13
% BEGIN ROW 14
\begin{align*}
R_{14}={}&6-19x_{9}-35x_{9}^{2}-80x_{9}x_{11}+4x_{9}x_{24}-24x_{9}x_{25}+33x_{10}+35x_{10}^{2}+80x_{10}x_{11}+24x_{10}x_{24}\\
{}&-4x_{10}x_{25}+16x_{11}+32x_{11}x_{24}-32x_{11}x_{25}+12x_{24}+4x_{24}^{2}-8x_{25}-4x_{25}^{2}+x_{30}+4x_{30}^{2}-x_{31}
\end{align*}
% END ROW 14
% BEGIN ROW 15
\begin{align*}
R_{15}={}&90-10x_{9}-2x_{9}^{2}-24x_{9}x_{10}-32x_{9}x_{11}-8x_{9}x_{24}-72x_{10}-18x_{10}^{2}-84x_{10}x_{11}-48x_{10}x_{24}-96x_{11}\\
{}&-74x_{11}^{2}-64x_{11}x_{24}-78x_{24}-8x_{24}^{2}-4x_{30}^{2}+x_{31}
\end{align*}
% END ROW 15
% BEGIN ROW 16
\begin{align*}
R_{16}={}&-12+16x_{9}-4x_{9}^{2}-8x_{9}x_{11}-8x_{9}x_{20}-8x_{9}x_{24}+16x_{11}-4x_{11}^{2}-8x_{11}x_{20}-8x_{11}x_{24}+16x_{20}\\
{}&-4x_{20}^{2}-8x_{20}x_{24}+16x_{24}-4x_{24}^{2}+x_{32}
\end{align*}
% END ROW 16
% BEGIN ROW 17
\begin{align*}
R_{17}={}&25-56x_{9}-48x_{9}^{2}-128x_{9}x_{11}-64x_{9}x_{24}-32x_{9}x_{25}-32x_{9}x_{32}-90x_{11}-52x_{11}^{2}-48x_{11}x_{24}\\
{}&-8x_{11}x_{25}-40x_{11}x_{32}-24x_{24}-8x_{24}^{2}-16x_{24}x_{32}-38x_{25}+4x_{25}^{2}-8x_{25}x_{32}-2x_{28}-4x_{32}-4x_{32}^{2}\\
{}&+x_{33}
\end{align*}
% END ROW 17
% BEGIN ROW 18
\begin{align*}
R_{18}={}&-18x_{9}-8x_{9}^{2}-64x_{9}x_{10}+16x_{9}x_{20}-32x_{10}-64x_{10}^{2}-2x_{20}+8x_{20}^{2}+2x_{21}-x_{37}
\end{align*}
% END ROW 18
% BEGIN ROW 19
\begin{align*}
R_{19}={}&-14-14x_{9}+23x_{9}^{2}+32x_{9}x_{10}+8x_{9}x_{20}+44x_{9}x_{21}-32x_{10}+64x_{10}x_{21}-8x_{20}+16x_{20}x_{21}+16x_{21}\\
{}&-4x_{21}^{2}+x_{34}
\end{align*}
% END ROW 19
% BEGIN ROW 20
\begin{align*}
R_{20}={}&1007-26x_{9}-4x_{9}^{2}-46x_{9}x_{10}-256x_{9}x_{20}+284x_{9}x_{21}+360x_{9}x_{34}+558x_{10}+41x_{10}^{2}\\
{}&-224x_{10}x_{20}+320x_{10}x_{21}+512x_{10}x_{34}-105x_{20}-28x_{20}^{2}+80x_{20}x_{21}+128x_{20}x_{34}-153x_{21}-20x_{21}^{2}\\
{}&-64x_{21}x_{34}+71x_{34}-4x_{34}^{2}+x_{35}
\end{align*}
% END ROW 20
% BEGIN ROW 21
\begin{align*}
R_{21}={}&-x_{10}^{2}-4x_{10}x_{35}-4x_{35}^{2}+x_{36}
\end{align*}
% END ROW 21
% BEGIN ROW 22
\begin{align*}
R_{22}={}&-3-6x_{9}-28x_{9}^{2}-32x_{9}x_{10}-40x_{9}x_{20}-32x_{10}x_{20}-6x_{20}-12x_{20}^{2}-x_{21}+x_{37}
\end{align*}
% END ROW 22
% BEGIN ROW 23
\begin{align*}
R_{23}={}&49-34x_{10}^{2}-44x_{10}x_{11}-24x_{10}x_{24}-32x_{10}x_{37}+6x_{11}-10x_{11}^{2}-24x_{11}x_{24}-24x_{11}x_{37}-32x_{24}\\
{}&-8x_{24}x_{37}+4x_{37}-4x_{37}^{2}+x_{38}-x_{44}
\end{align*}
% END ROW 23
% BEGIN ROW 24
\begin{align*}
R_{24}={}&8x_{9}-64x_{9}^{2}-96x_{9}x_{11}-32x_{9}x_{25}-32x_{9}x_{38}+6x_{11}-36x_{11}^{2}-24x_{11}x_{25}-24x_{11}x_{38}+2x_{25}\\
{}&-4x_{25}^{2}-8x_{25}x_{38}+4x_{38}-4x_{38}^{2}+x_{39}
\end{align*}
% END ROW 24
% BEGIN ROW 25
\begin{align*}
R_{25}={}&-x_{11}^{2}-2x_{11}x_{12}-x_{12}^{2}+x_{40}
\end{align*}
% END ROW 25
% BEGIN ROW 26
\begin{align*}
R_{26}={}&-11+3x_{9}^{2}+8x_{9}x_{11}-8x_{9}x_{25}-4x_{9}x_{40}+4x_{10}^{2}+8x_{10}x_{24}-8x_{10}x_{25}-8x_{11}+2x_{11}^{2}-8x_{11}x_{40}\\
{}&+4x_{24}^{2}-8x_{24}x_{25}+18x_{25}-4x_{25}^{2}+4x_{26}^{2}-x_{27}+4x_{40}-4x_{40}^{2}+x_{41}
\end{align*}
% END ROW 26
% BEGIN ROW 27
\begin{align*}
R_{27}={}&4x_{9}-18x_{9}^{2}-40x_{9}x_{11}-32x_{9}x_{25}-40x_{9}x_{40}-16x_{9}x_{41}+4x_{11}-24x_{11}^{2}-32x_{11}x_{25}-48x_{11}x_{40}\\
{}&-16x_{11}x_{41}-16x_{25}^{2}-32x_{25}x_{40}-16x_{25}x_{41}+4x_{40}-24x_{40}^{2}-16x_{40}x_{41}-4x_{41}^{2}+x_{42}
\end{align*}
% END ROW 27
% BEGIN ROW 28
\begin{align*}
R_{28}={}&x_{11}-25x_{11}^{2}-40x_{11}x_{25}-20x_{11}x_{28}-16x_{25}^{2}-16x_{25}x_{28}-4x_{28}^{2}+x_{43}
\end{align*}
% END ROW 28
% BEGIN ROW 29
\begin{align*}
R_{29}={}&-1+8x_{10}-4x_{10}^{2}-8x_{10}x_{24}+8x_{24}-4x_{24}^{2}+x_{44}
\end{align*}
% END ROW 29
% BEGIN ROW 30
\begin{align*}
R_{30}={}&-1-16x_{10}^{2}-8x_{10}x_{11}-32x_{10}x_{24}-16x_{10}x_{44}-x_{11}^{2}-8x_{11}x_{24}-4x_{11}x_{44}-16x_{24}^{2}-16x_{24}x_{44}\\
{}&-4x_{44}^{2}+x_{45}
\end{align*}
% END ROW 30
% BEGIN ROW 31
\begin{align*}
R_{31}={}&3+x_{10}-2x_{10}x_{11}-4x_{10}x_{45}-3x_{11}^{2}-8x_{11}x_{45}-x_{24}-4x_{45}^{2}+x_{46}
\end{align*}
% END ROW 31
% BEGIN ROW 32
\begin{align*}
R_{32}={}&-1-4x_{46}-4x_{46}^{2}+x_{47}
\end{align*}
% END ROW 32
% BEGIN ROW 33
\begin{align*}
R_{33}={}&-1-x_{11}^{2}-8x_{11}x_{46}-4x_{11}x_{47}-16x_{46}^{2}-16x_{46}x_{47}-4x_{47}^{2}+x_{48}
\end{align*}
% END ROW 33
% BEGIN ROW 34
\begin{align*}
R_{34}={}&4+x_{10}-2x_{10}x_{11}-4x_{10}x_{48}-3x_{11}^{2}-8x_{11}x_{48}-x_{24}-4x_{48}^{2}+x_{49}
\end{align*}
% END ROW 34
% BEGIN ROW 35
\begin{align*}
R_{35}={}&-4x_{43}^{2}-8x_{43}x_{49}-4x_{49}^{2}+x_{50}
\end{align*}
% END ROW 35
% BEGIN ROW 36
\begin{align*}
R_{36}={}&-9+7x_{9}-x_{9}^{2}-8x_{9}x_{43}-4x_{9}x_{50}+24x_{43}-16x_{43}^{2}-16x_{43}x_{50}+12x_{50}-4x_{50}^{2}+x_{51}
\end{align*}
% END ROW 36
% BEGIN ROW 37
\begin{align*}
R_{37}={}&-1-4x_{51}-4x_{51}^{2}+x_{52}
\end{align*}
% END ROW 37
% BEGIN ROW 38
\begin{align*}
R_{38}={}&-x_{11}^{2}-8x_{11}x_{51}-4x_{11}x_{52}-16x_{51}^{2}-16x_{51}x_{52}-4x_{52}^{2}+x_{53}
\end{align*}
% END ROW 38
% BEGIN ROW 39
\begin{align*}
R_{39}={}&-x_{9}-x_{9}^{2}-4x_{9}x_{11}-4x_{9}x_{53}-4x_{11}-4x_{11}^{2}-8x_{11}x_{53}-4x_{53}-4x_{53}^{2}+x_{54}
\end{align*}
% END ROW 39
% BEGIN ROW 40
\begin{align*}
R_{40}={}&-10-12x_{54}-4x_{54}^{2}+x_{55}
\end{align*}
% END ROW 40
% BEGIN ROW 41
\begin{align*}
R_{41}={}&2-4x_{10}-16x_{10}x_{54}-8x_{10}x_{55}-4x_{24}-16x_{24}x_{54}-8x_{24}x_{55}-x_{44}+8x_{54}-16x_{54}^{2}-16x_{54}x_{55}+4x_{55}\\
{}&-4x_{55}^{2}+x_{56}
\end{align*}
% END ROW 41
% BEGIN ROW 42
\begin{align*}
R_{42}={}&-10+x_{10}-25x_{10}^{2}-40x_{10}x_{24}-20x_{10}x_{56}-16x_{24}^{2}-16x_{24}x_{56}-12x_{49}-4x_{49}^{2}-4x_{56}^{2}+x_{57}
\end{align*}
% END ROW 42
% BEGIN ROW 43
\begin{align*}
R_{43}={}&21+24x_{49}-8x_{49}^{2}-16x_{49}x_{57}-2x_{57}-4x_{57}^{2}+x_{58}
\end{align*}
% END ROW 43
% BEGIN ROW 44
\begin{align*}
R_{44}={}&1+4x_{9}-4x_{9}^{2}-8x_{9}x_{20}-8x_{9}x_{60}+x_{11}-x_{11}^{2}-4x_{11}x_{58}+4x_{20}-4x_{20}^{2}-8x_{20}x_{60}-4x_{58}^{2}+x_{59}\\
{}&+4x_{60}-4x_{60}^{2}
\end{align*}
% END ROW 44
% BEGIN ROW 45
\begin{align*}
R_{45}={}&-4x_{9}-4x_{9}^{2}-4x_{9}x_{11}-8x_{9}x_{20}-x_{11}-x_{11}^{2}-4x_{11}x_{20}-4x_{20}-4x_{20}^{2}+x_{60}
\end{align*}
% END ROW 45
% BEGIN ROW 46
\begin{align*}
R_{46}={}&4x_{9}-24x_{9}^{2}-48x_{9}x_{20}-16x_{9}x_{59}-16x_{9}x_{60}+4x_{20}-24x_{20}^{2}-16x_{20}x_{59}-16x_{20}x_{60}-4x_{59}^{2}\\
{}&+8x_{60}-8x_{60}^{2}+x_{61}
\end{align*}
% END ROW 46
% BEGIN ROW 47
\begin{align*}
R_{47}={}&-2-4x_{61}-4x_{61}^{2}+x_{62}
\end{align*}
% END ROW 47
% BEGIN ROW 48
\begin{align*}
R_{48}={}&-x_{2}x_{14}+x_{63}
\end{align*}
% END ROW 48
% BEGIN ROW 49
\begin{align*}
R_{49}={}&x_{13}-x_{14}-2x_{23}x_{63}-x_{64}
\end{align*}
% END ROW 49
% BEGIN ROW 50
\begin{align*}
R_{50}={}&x_{13}-x_{14}-2x_{27}x_{63}-x_{65}
\end{align*}
% END ROW 50
% BEGIN ROW 51
\begin{align*}
R_{51}={}&-2x_{29}x_{63}+2x_{42}x_{63}-2x_{63}-x_{66}+x_{71}
\end{align*}
% END ROW 51
% BEGIN ROW 52
\begin{align*}
R_{52}={}&x_{13}-x_{14}-2x_{31}x_{63}-x_{67}
\end{align*}
% END ROW 52
% BEGIN ROW 53
\begin{align*}
R_{53}={}&x_{13}-x_{14}-2x_{33}x_{63}+6x_{63}-x_{68}
\end{align*}
% END ROW 53
% BEGIN ROW 54
\begin{align*}
R_{54}={}&x_{13}-x_{14}-4x_{35}x_{63}-2x_{36}x_{63}-x_{69}
\end{align*}
% END ROW 54
% BEGIN ROW 55
\begin{align*}
R_{55}={}&x_{13}-x_{14}-2x_{39}x_{63}-x_{70}
\end{align*}
% END ROW 55
% BEGIN ROW 56
\begin{align*}
R_{56}={}&x_{13}-x_{14}-2x_{42}x_{63}-4x_{63}-x_{71}
\end{align*}
% END ROW 56
% BEGIN ROW 57
\begin{align*}
R_{57}={}&-x_{9}x_{10}+x_{9}x_{11}+x_{10}+x_{10}x_{11}-x_{11}-x_{11}^{2}+x_{72}
\end{align*}
% END ROW 57
% BEGIN ROW 58
\begin{align*}
R_{58}={}&1+2x_{12}+x_{12}^{2}-x_{72}^{2}+x_{73}
\end{align*}
% END ROW 58
% BEGIN ROW 59
\begin{align*}
R_{59}={}&-x_{13}^{2}+2x_{13}x_{14}+2x_{13}x_{75}-x_{14}^{2}-2x_{14}x_{75}-x_{73}^{2}+x_{74}-x_{75}^{2}
\end{align*}
% END ROW 59
% BEGIN ROW 60
\begin{align*}
R_{60}={}&-4x_{61}x_{63}-2x_{62}x_{63}+2x_{63}+x_{75}
\end{align*}
% END ROW 60
% BEGIN ROW 61
\begin{align*}
R_{61}={}&-x_{10}x_{63}-x_{14}+x_{76}x_{77}+x_{78}
\end{align*}
% END ROW 61
% BEGIN ROW 62
\begin{align*}
R_{62}={}&-1+2x_{2}-2x_{2}x_{9}-2x_{2}x_{20}+x_{76}
\end{align*}
% END ROW 62
% BEGIN ROW 63
\begin{align*}
R_{63}={}&-x_{9}x_{213}-x_{13}+x_{77}+x_{213}
\end{align*}
% END ROW 63
% BEGIN ROW 64
\begin{align*}
R_{64}={}&-6x_{9}x_{63}-8x_{10}x_{63}+x_{10}x_{213}+x_{13}-x_{14}-2x_{20}x_{63}+x_{79}
\end{align*}
% END ROW 64
% BEGIN ROW 65
\begin{align*}
R_{65}={}&-x_{13}x_{64}+x_{14}x_{64}+x_{80}
\end{align*}
% END ROW 65
% BEGIN ROW 66
\begin{align*}
R_{66}={}&-x_{65}x_{80}+x_{81}
\end{align*}
% END ROW 66
% BEGIN ROW 67
\begin{align*}
R_{67}={}&-x_{66}x_{81}+x_{82}
\end{align*}
% END ROW 67
% BEGIN ROW 68
\begin{align*}
R_{68}={}&-x_{67}x_{82}+x_{83}
\end{align*}
% END ROW 68
% BEGIN ROW 69
\begin{align*}
R_{69}={}&-x_{68}x_{83}+x_{84}
\end{align*}
% END ROW 69
% BEGIN ROW 70
\begin{align*}
R_{70}={}&-x_{69}x_{84}+x_{85}
\end{align*}
% END ROW 70
% BEGIN ROW 71
\begin{align*}
R_{71}={}&-x_{70}x_{85}+x_{86}
\end{align*}
% END ROW 71
% BEGIN ROW 72
\begin{align*}
R_{72}={}&-x_{71}x_{86}+x_{87}
\end{align*}
% END ROW 72
% BEGIN ROW 73
\begin{align*}
R_{73}={}&-x_{78}x_{87}+x_{88}
\end{align*}
% END ROW 73
% BEGIN ROW 74
\begin{align*}
R_{74}={}&-x_{79}x_{88}+x_{89}
\end{align*}
% END ROW 74
% BEGIN ROW 75
\begin{align*}
R_{75}={}&-x_{89}^{2}+x_{90}
\end{align*}
% END ROW 75
% BEGIN ROW 76
\begin{align*}
R_{76}={}&-x_{5}-x_{5}x_{17}+x_{17}+x_{91}
\end{align*}
% END ROW 76
% BEGIN ROW 77
\begin{align*}
R_{77}={}&-x_{15}x_{91}+x_{92}
\end{align*}
% END ROW 77
% BEGIN ROW 78
\begin{align*}
R_{78}={}&-2x_{1}^{2}+2x_{1}x_{13}+2x_{1}x_{14}+2x_{1}x_{92}+2x_{1}x_{93}-x_{13}^{2}-2x_{13}x_{92}-x_{14}^{2}-2x_{14}x_{93}+x_{18}x_{94}-x_{74}x_{90}\\
{}&-x_{92}^{2}-x_{93}^{2}
\end{align*}
% END ROW 78
% BEGIN ROW 79
\begin{align*}
R_{79}={}&-x_{5}x_{16}+x_{16}-x_{16}x_{91}+x_{93}
\end{align*}
% END ROW 79
% BEGIN ROW 80
\begin{align*}
R_{80}={}&-x_{95}x_{98}+x_{98}+x_{100}
\end{align*}
% END ROW 80
% BEGIN ROW 81
\begin{align*}
R_{81}={}&-x_{4}x_{121}+x_{96}x_{203}-x_{99}-x_{100}x_{121}
\end{align*}
% END ROW 81
% BEGIN ROW 82
\begin{align*}
R_{82}={}&x_{101}x_{203}-x_{102}x_{124}-x_{103}
\end{align*}
% END ROW 82
% BEGIN ROW 83
\begin{align*}
R_{83}={}&x_{104}x_{203}-x_{105}x_{124}-x_{106}
\end{align*}
% END ROW 83
% BEGIN ROW 84
\begin{align*}
R_{84}={}&x_{107}x_{203}-x_{108}x_{124}-x_{109}
\end{align*}
% END ROW 84
% BEGIN ROW 85
\begin{align*}
R_{85}={}&x_{110}x_{203}-x_{111}x_{124}-x_{112}
\end{align*}
% END ROW 85
% BEGIN ROW 86
\begin{align*}
R_{86}={}&x_{113}x_{204}-x_{115}x_{124}-x_{116}
\end{align*}
% END ROW 86
% BEGIN ROW 87
\begin{align*}
R_{87}={}&x_{117}x_{204}-x_{118}x_{124}-x_{119}
\end{align*}
% END ROW 87
% BEGIN ROW 88
\begin{align*}
R_{88}={}&-x_{18}x_{121}+x_{124}
\end{align*}
% END ROW 88
% BEGIN ROW 89
\begin{align*}
R_{89}={}&x_{120}-x_{122}x_{124}-x_{123}
\end{align*}
% END ROW 89
% BEGIN ROW 90
\begin{align*}
R_{90}={}&-x_{124}x_{126}+x_{125}-x_{127}
\end{align*}
% END ROW 90
% BEGIN ROW 91
\begin{align*}
R_{91}={}&-x_{124}x_{129}+x_{128}x_{203}-x_{130}
\end{align*}
% END ROW 91
% BEGIN ROW 92
\begin{align*}
R_{92}={}&-x_{124}x_{132}+x_{131}x_{203}-x_{133}
\end{align*}
% END ROW 92
% BEGIN ROW 93
\begin{align*}
R_{93}={}&x_{134}-2x_{134}x_{136}+x_{137}
\end{align*}
% END ROW 93
% BEGIN ROW 94
\begin{align*}
R_{94}={}&-x_{134}x_{137}-x_{137}+x_{138}
\end{align*}
% END ROW 94
% BEGIN ROW 95
\begin{align*}
R_{95}={}&2x_{135}+x_{135}^{2}+x_{137}^{2}-x_{138}^{2}
\end{align*}
% END ROW 95
% BEGIN ROW 96
\begin{align*}
R_{96}={}&-x_{135}x_{140}-x_{140}+x_{141}
\end{align*}
% END ROW 96
% BEGIN ROW 97
\begin{align*}
R_{97}={}&-1+x_{139}^{2}+2x_{139}x_{141}+x_{140}^{2}
\end{align*}
% END ROW 97
% BEGIN ROW 98
\begin{align*}
R_{98}={}&x_{3}^{2}-2x_{3}x_{135}-2x_{135}+x_{142}
\end{align*}
% END ROW 98
% BEGIN ROW 99
\begin{align*}
R_{99}={}&-x_{3}x_{140}+x_{5}-x_{139}-x_{140}+x_{142}x_{143}
\end{align*}
% END ROW 99
% BEGIN ROW 100
\begin{align*}
R_{100}={}&1-2x_{4}+x_{4}^{2}-2x_{4}x_{135}+x_{144}
\end{align*}
% END ROW 100
% BEGIN ROW 101
\begin{align*}
R_{101}={}&-x_{4}x_{176}+x_{18}+x_{144}x_{145}-x_{175}
\end{align*}
% END ROW 101
% BEGIN ROW 102
\begin{align*}
R_{102}={}&-1+2x_{135}x_{187}-x_{146}+2x_{187}-x_{187}^{2}
\end{align*}
% END ROW 102
% BEGIN ROW 103
\begin{align*}
R_{103}={}&-1+x_{95}+x_{146}x_{147}-x_{175}-x_{176}x_{187}
\end{align*}
% END ROW 103
% BEGIN ROW 104
\begin{align*}
R_{104}={}&-x_{135}x_{149}-x_{149}+x_{150}
\end{align*}
% END ROW 104
% BEGIN ROW 105
\begin{align*}
R_{105}={}&-1+x_{148}^{2}+2x_{148}x_{150}+x_{149}^{2}
\end{align*}
% END ROW 105
% BEGIN ROW 106
\begin{align*}
R_{106}={}&1+x_{135}-x_{135}x_{152}-x_{152}+x_{153}
\end{align*}
% END ROW 106
% BEGIN ROW 107
\begin{align*}
R_{107}={}&x_{151}^{2}+2x_{151}x_{153}-2x_{152}+x_{152}^{2}
\end{align*}
% END ROW 107
% BEGIN ROW 108
\begin{align*}
R_{108}={}&-x_{135}x_{155}-x_{155}+x_{156}
\end{align*}
% END ROW 108
% BEGIN ROW 109
\begin{align*}
R_{109}={}&-1+x_{154}^{2}+2x_{154}x_{156}+x_{155}^{2}
\end{align*}
% END ROW 109
% BEGIN ROW 110
\begin{align*}
R_{110}={}&-x_{135}x_{158}-x_{158}+x_{159}
\end{align*}
% END ROW 110
% BEGIN ROW 111
\begin{align*}
R_{111}={}&-1+x_{157}^{2}+2x_{157}x_{159}+x_{158}^{2}
\end{align*}
% END ROW 111
% BEGIN ROW 112
\begin{align*}
R_{112}={}&-x_{135}x_{161}-x_{161}+x_{162}
\end{align*}
% END ROW 112
% BEGIN ROW 113
\begin{align*}
R_{113}={}&-1+x_{160}^{2}+2x_{160}x_{162}+x_{161}^{2}
\end{align*}
% END ROW 113
% BEGIN ROW 114
\begin{align*}
R_{114}={}&-x_{135}x_{164}-x_{164}+x_{165}
\end{align*}
% END ROW 114
% BEGIN ROW 115
\begin{align*}
R_{115}={}&-1+x_{163}^{2}+2x_{163}x_{165}+x_{164}^{2}
\end{align*}
% END ROW 115
% BEGIN ROW 116
\begin{align*}
R_{116}={}&-x_{135}x_{167}-x_{167}+x_{168}
\end{align*}
% END ROW 116
% BEGIN ROW 117
\begin{align*}
R_{117}={}&-1+x_{166}^{2}+2x_{166}x_{168}+x_{167}^{2}
\end{align*}
% END ROW 117
% BEGIN ROW 118
\begin{align*}
R_{118}={}&-x_{135}x_{170}-x_{170}+x_{171}
\end{align*}
% END ROW 118
% BEGIN ROW 119
\begin{align*}
R_{119}={}&-1+x_{169}^{2}+2x_{169}x_{171}+x_{170}^{2}
\end{align*}
% END ROW 119
% BEGIN ROW 120
\begin{align*}
R_{120}={}&-x_{135}x_{173}-x_{173}+x_{174}
\end{align*}
% END ROW 120
% BEGIN ROW 121
\begin{align*}
R_{121}={}&-1+x_{172}^{2}+2x_{172}x_{174}+x_{173}^{2}
\end{align*}
% END ROW 121
% BEGIN ROW 122
\begin{align*}
R_{122}={}&1-x_{135}x_{176}-x_{176}+x_{177}
\end{align*}
% END ROW 122
% BEGIN ROW 123
\begin{align*}
R_{123}={}&-1+2x_{175}+x_{175}^{2}+2x_{175}x_{177}+x_{176}^{2}
\end{align*}
% END ROW 123
% BEGIN ROW 124
\begin{align*}
R_{124}={}&-x_{135}x_{179}-x_{179}+x_{180}
\end{align*}
% END ROW 124
% BEGIN ROW 125
\begin{align*}
R_{125}={}&-1+x_{178}^{2}+2x_{178}x_{180}+x_{179}^{2}
\end{align*}
% END ROW 125
% BEGIN ROW 126
\begin{align*}
R_{126}={}&-x_{135}x_{182}-x_{182}+x_{183}
\end{align*}
% END ROW 126
% BEGIN ROW 127
\begin{align*}
R_{127}={}&-1+x_{181}^{2}+2x_{181}x_{183}+x_{182}^{2}
\end{align*}
% END ROW 127
% BEGIN ROW 128
\begin{align*}
R_{128}={}&-x_{135}x_{185}-x_{185}+x_{186}
\end{align*}
% END ROW 128
% BEGIN ROW 129
\begin{align*}
R_{129}={}&-1+x_{184}^{2}+2x_{184}x_{186}+x_{185}^{2}
\end{align*}
% END ROW 129
% BEGIN ROW 130
\begin{align*}
R_{130}={}&-2-2x_{0}-3x_{0}^{2}-2x_{0}x_{3}-x_{1}^{2}-2x_{3}-2x_{3}^{2}-x_{4}^{2}-2x_{5}^{2}+2x_{9}-x_{9}^{2}+2x_{10}x_{11}-x_{12}^{2}-x_{13}^{2}\\
{}&-x_{14}^{2}-x_{15}^{2}-x_{16}^{2}-x_{17}^{2}-x_{18}^{2}-x_{24}+2x_{95}-2x_{95}^{2}-x_{96}^{2}-x_{97}^{2}-x_{101}^{2}-x_{104}^{2}-x_{107}^{2}-x_{110}^{2}\\
{}&-x_{113}^{2}-x_{114}^{2}-x_{117}^{2}-x_{120}^{2}-x_{121}^{2}-x_{125}^{2}-x_{128}^{2}-x_{131}^{2}+x_{134}+x_{140}^{2}+x_{149}^{2}-2x_{152}+x_{152}^{2}+x_{155}^{2}\\
{}&+x_{158}^{2}+x_{161}^{2}+x_{164}^{2}+x_{167}^{2}+x_{170}^{2}+x_{173}^{2}+x_{176}^{2}+x_{179}^{2}+x_{182}^{2}-x_{185}^{2}-x_{187}^{2}+x_{205}^{2}+x_{206}^{2}+x_{207}\\
{}&+x_{207}^{2}
\end{align*}
% END ROW 130
% BEGIN ROW 131
\begin{align*}
R_{131}={}&x_{4}-2x_{4}x_{9}-2x_{4}x_{10}-2x_{4}x_{11}-2x_{4}x_{12}-2x_{4}x_{13}-2x_{4}x_{14}-2x_{4}x_{15}-2x_{4}x_{16}-2x_{4}x_{17}+x_{187}
\end{align*}
% END ROW 131
% BEGIN ROW 132
\begin{align*}
R_{132}={}&-16+x_{0}-x_{188}^{2}
\end{align*}
% END ROW 132
% BEGIN ROW 133
\begin{align*}
R_{133}={}&-x_{135}x_{190}-x_{190}+x_{199}
\end{align*}
% END ROW 133
% BEGIN ROW 134
\begin{align*}
R_{134}={}&-1+x_{189}^{2}+x_{190}^{2}-x_{199}^{2}
\end{align*}
% END ROW 134
% BEGIN ROW 135
\begin{align*}
R_{135}={}&-x_{192}-x_{192}x_{193}+x_{200}
\end{align*}
% END ROW 135
% BEGIN ROW 136
\begin{align*}
R_{136}={}&-1+x_{191}^{2}+x_{192}^{2}-x_{200}^{2}
\end{align*}
% END ROW 136
% BEGIN ROW 137
\begin{align*}
R_{137}={}&-x_{185}^{2}+x_{201}
\end{align*}
% END ROW 137
% BEGIN ROW 138
\begin{align*}
R_{138}={}&x_{190}-2x_{195}x_{201}-x_{201}
\end{align*}
% END ROW 138
% BEGIN ROW 139
\begin{align*}
R_{139}={}&-4x_{190}-8x_{190}x_{194}+x_{193}
\end{align*}
% END ROW 139
% BEGIN ROW 140
\begin{align*}
R_{140}={}&-x_{135}-x_{189}x_{196}+x_{193}
\end{align*}
% END ROW 140
% BEGIN ROW 141
\begin{align*}
R_{141}={}&-x_{184}-x_{186}-x_{189}x_{197}+x_{191}
\end{align*}
% END ROW 141
% BEGIN ROW 142
\begin{align*}
R_{142}={}&-9x_{9}-9x_{10}-4x_{20}-2x_{23}-2x_{27}-2x_{29}-2x_{31}-2x_{33}-4x_{35}-2x_{36}-2x_{39}-2x_{42}-4x_{61}-2x_{62}\\
{}&-4x_{185}x_{198}+x_{192}
\end{align*}
% END ROW 142
% BEGIN ROW 143
\begin{align*}
R_{143}={}&20-2x_{0}-x_{0}^{2}-2x_{0}x_{140}-2x_{0}x_{149}-2x_{1}x_{176}-2x_{3}-x_{3}^{2}-2x_{3}x_{149}-x_{4}^{2}-x_{5}^{2}-2x_{5}x_{167}+2x_{9}\\
{}&-2x_{9}x_{152}-18x_{9}x_{185}-2x_{10}x_{155}-18x_{10}x_{185}-2x_{11}x_{158}-2x_{12}x_{161}-2x_{13}x_{164}-2x_{14}x_{170}-2x_{15}x_{173}\\
{}&-2x_{16}x_{179}-2x_{17}x_{182}-x_{18}^{2}-8x_{20}x_{185}-4x_{23}x_{185}-4x_{27}x_{185}-4x_{29}x_{185}-4x_{31}x_{185}-4x_{33}x_{185}\\
{}&-8x_{35}x_{185}-4x_{36}x_{185}-4x_{39}x_{185}-4x_{42}x_{185}-8x_{61}x_{185}-4x_{62}x_{185}+2x_{95}-2x_{95}^{2}-x_{96}^{2}-x_{97}^{2}-x_{101}^{2}\\
{}&-x_{104}^{2}-x_{107}^{2}-x_{110}^{2}-x_{113}^{2}-x_{114}^{2}-x_{117}^{2}-x_{120}^{2}-x_{121}^{2}-x_{125}^{2}-x_{128}^{2}-x_{131}^{2}+x_{134}-x_{135}x_{202}\\
{}&-2x_{139}-2x_{141}+x_{141}^{2}-2x_{148}-2x_{150}+x_{150}^{2}-2x_{151}+2x_{152}-2x_{153}+x_{153}^{2}-2x_{154}-2x_{156}+x_{156}^{2}\\
{}&-2x_{157}-2x_{159}+x_{159}^{2}-2x_{160}-2x_{162}+x_{162}^{2}-2x_{163}-2x_{165}+x_{165}^{2}-2x_{166}-2x_{168}+x_{168}^{2}-2x_{169}\\
{}&-2x_{171}+x_{171}^{2}-2x_{172}-2x_{174}+x_{174}^{2}-2x_{175}+x_{177}^{2}-2x_{178}-2x_{180}+x_{180}^{2}-2x_{181}-2x_{183}+x_{183}^{2}\\
{}&-2x_{184}-2x_{186}+x_{186}^{2}-x_{187}^{2}
\end{align*}
% END ROW 143
% BEGIN ROW 144
\begin{align*}
R_{144}={}&-x_{97}x_{203}+x_{121}
\end{align*}
% END ROW 144
% BEGIN ROW 145
\begin{align*}
R_{145}={}&-x_{114}x_{204}+x_{121}
\end{align*}
% END ROW 145
% BEGIN ROW 146
\begin{align*}
R_{146}={}&-2-81x_{9}^{2}-162x_{9}x_{10}-72x_{9}x_{20}-36x_{9}x_{23}-36x_{9}x_{27}-36x_{9}x_{29}-36x_{9}x_{31}-36x_{9}x_{33}-72x_{9}x_{35}\\
{}&-36x_{9}x_{36}-36x_{9}x_{39}-36x_{9}x_{42}-72x_{9}x_{61}-36x_{9}x_{62}-81x_{10}^{2}-72x_{10}x_{20}-36x_{10}x_{23}-36x_{10}x_{27}\\
{}&-36x_{10}x_{29}-36x_{10}x_{31}-36x_{10}x_{33}-72x_{10}x_{35}-36x_{10}x_{36}-36x_{10}x_{39}-36x_{10}x_{42}-72x_{10}x_{61}\\
{}&-36x_{10}x_{62}-16x_{20}^{2}-16x_{20}x_{23}-16x_{20}x_{27}-16x_{20}x_{29}-16x_{20}x_{31}-16x_{20}x_{33}-32x_{20}x_{35}-16x_{20}x_{36}\\
{}&-16x_{20}x_{39}-16x_{20}x_{42}-32x_{20}x_{61}-16x_{20}x_{62}-4x_{23}^{2}-8x_{23}x_{27}-8x_{23}x_{29}-8x_{23}x_{31}-8x_{23}x_{33}\\
{}&-16x_{23}x_{35}-8x_{23}x_{36}-8x_{23}x_{39}-8x_{23}x_{42}-16x_{23}x_{61}-8x_{23}x_{62}-4x_{27}^{2}-8x_{27}x_{29}-8x_{27}x_{31}\\
{}&-8x_{27}x_{33}-16x_{27}x_{35}-8x_{27}x_{36}-8x_{27}x_{39}-8x_{27}x_{42}-16x_{27}x_{61}-8x_{27}x_{62}-4x_{29}^{2}-8x_{29}x_{31}\\
{}&-8x_{29}x_{33}-16x_{29}x_{35}-8x_{29}x_{36}-8x_{29}x_{39}-8x_{29}x_{42}-16x_{29}x_{61}-8x_{29}x_{62}-4x_{31}^{2}-8x_{31}x_{33}\\
{}&-16x_{31}x_{35}-8x_{31}x_{36}-8x_{31}x_{39}-8x_{31}x_{42}-16x_{31}x_{61}-8x_{31}x_{62}-4x_{33}^{2}-16x_{33}x_{35}-8x_{33}x_{36}\\
{}&-8x_{33}x_{39}-8x_{33}x_{42}-16x_{33}x_{61}-8x_{33}x_{62}-16x_{35}^{2}-16x_{35}x_{36}-16x_{35}x_{39}-16x_{35}x_{42}-32x_{35}x_{61}\\
{}&-16x_{35}x_{62}-4x_{36}^{2}-8x_{36}x_{39}-8x_{36}x_{42}-16x_{36}x_{61}-8x_{36}x_{62}-4x_{39}^{2}-8x_{39}x_{42}-16x_{39}x_{61}\\
{}&-8x_{39}x_{62}-4x_{42}^{2}-16x_{42}x_{61}-8x_{42}x_{62}-16x_{61}^{2}-16x_{61}x_{62}-4x_{62}^{2}-x_{140}^{2}-x_{149}^{2}+2x_{152}-x_{152}^{2}\\
{}&-x_{155}^{2}-x_{158}^{2}-x_{161}^{2}-x_{164}^{2}-x_{167}^{2}-x_{170}^{2}-x_{173}^{2}-x_{176}^{2}-x_{179}^{2}-x_{182}^{2}+x_{185}^{2}-x_{205}^{2}-x_{206}^{2}-x_{207}\\
{}&-x_{207}^{2}
\end{align*}
% END ROW 146
% BEGIN ROW 147
\begin{align*}
R_{147}={}&-1-4x_{99}^{2}-4x_{103}^{2}-4x_{106}^{2}-4x_{109}^{2}-4x_{112}^{2}-4x_{116}^{2}-4x_{119}^{2}+x_{121}^{2}-4x_{123}^{2}-4x_{127}^{2}-4x_{130}^{2}\\
{}&-4x_{133}^{2}-x_{208}^{2}-x_{209}^{2}-x_{210}-x_{210}^{2}-4x_{212}^{2}
\end{align*}
% END ROW 147
% BEGIN ROW 148
\begin{align*}
R_{148}={}&-x_{4}x_{211}+x_{121}
\end{align*}
% END ROW 148
% BEGIN ROW 149
\begin{align*}
R_{149}={}&x_{2}x_{121}-x_{5}x_{211}-x_{212}
\end{align*}
% END ROW 149
% BEGIN ROW 150
\begin{align*}
R_{150}={}&-x_{2}x_{13}+x_{213}
\end{align*}
% END ROW 150
% BEGIN ROW 151
\begin{align*}
R_{151}={}&x_{95}x_{215}-2x_{96}x_{149}+2x_{101}-2x_{101}x_{152}-2x_{104}x_{155}-2x_{107}x_{158}-2x_{110}x_{161}-2x_{113}x_{164}-2x_{117}x_{170}\\
{}&-2x_{120}x_{173}-2x_{125}x_{179}-2x_{128}x_{182}-2x_{131}x_{185}-2x_{135}x_{215}+2x_{148}x_{149}+2x_{149}x_{150}-2x_{151}+2x_{151}x_{152}\\
{}&+2x_{152}x_{153}-2x_{153}+2x_{154}x_{155}+2x_{155}x_{156}+2x_{157}x_{158}+2x_{158}x_{159}+2x_{160}x_{161}+2x_{161}x_{162}+2x_{163}x_{164}\\
{}&+2x_{164}x_{165}+2x_{169}x_{170}+2x_{170}x_{171}+2x_{172}x_{173}+2x_{173}x_{174}+2x_{178}x_{179}+2x_{179}x_{180}+2x_{181}x_{182}+2x_{182}x_{183}\\
{}&+2x_{184}x_{185}+2x_{185}x_{186}+x_{214}-2x_{215}
\end{align*}
% END ROW 151
% BEGIN ROW 152
\begin{align*}
R_{152}={}&-1-x_{95}x_{214}+x_{96}^{2}-2x_{96}x_{148}+x_{101}^{2}-2x_{101}x_{151}+x_{104}^{2}-2x_{104}x_{154}+x_{107}^{2}-2x_{107}x_{157}+x_{110}^{2}\\
{}&-2x_{110}x_{160}+x_{113}^{2}-2x_{113}x_{163}+x_{117}^{2}-2x_{117}x_{169}+x_{120}^{2}-2x_{120}x_{172}+x_{125}^{2}-2x_{125}x_{178}+x_{128}^{2}\\
{}&-2x_{128}x_{181}+x_{131}^{2}-2x_{131}x_{184}+x_{148}^{2}-x_{149}^{2}+x_{151}^{2}+2x_{152}-x_{152}^{2}+x_{154}^{2}-x_{155}^{2}+x_{157}^{2}-x_{158}^{2}\\
{}&+x_{160}^{2}-x_{161}^{2}+x_{163}^{2}-x_{164}^{2}+x_{169}^{2}-x_{170}^{2}+x_{172}^{2}-x_{173}^{2}+x_{178}^{2}-x_{179}^{2}+x_{181}^{2}-x_{182}^{2}+x_{184}^{2}\\
{}&-x_{185}^{2}+x_{215}
\end{align*}
% END ROW 152
% BEGIN ROW 153
\begin{align*}
R_{153}={}&x_{95}x_{217}-2x_{97}x_{140}-2x_{114}x_{167}-2x_{121}x_{176}-2x_{135}x_{217}+2x_{139}x_{140}+2x_{140}x_{141}+2x_{166}x_{167}+2x_{167}x_{168}\\
{}&+2x_{175}x_{176}+2x_{176}+2x_{176}x_{177}+x_{216}-3x_{217}
\end{align*}
% END ROW 153
% BEGIN ROW 154
\begin{align*}
R_{154}={}&-x_{95}x_{216}+x_{97}^{2}-2x_{97}x_{139}+x_{114}^{2}-2x_{114}x_{166}+x_{121}^{2}-2x_{121}x_{175}+x_{139}^{2}-x_{140}^{2}+x_{166}^{2}-x_{167}^{2}+x_{175}^{2}\\
{}&-x_{176}^{2}+x_{216}+x_{217}
\end{align*}
% END ROW 154

\endgroup
\end{document}
